\documentclass[11pt,a4paper]{article}

% --- Standard Academic Packages ---
\usepackage[utf8]{inputenc}
\usepackage[margin=1in]{geometry}
\usepackage{amsmath,amssymb,amsfonts,amsthm}
\usepackage{booktabs}
\usepackage{graphicx}
\usepackage{algorithm}
\usepackage{algorithmic}
\usepackage{hyperref}
\usepackage{cite}
\usepackage{microtype}
\usepackage{subcaption}
\usepackage{xcolor}
\usepackage{array}
\usepackage{multirow}
\hyphenation{SIGMA-QGD}

% --- Theorem Environments ---
\newtheorem{theorem}{Theorem}
\newtheorem{lemma}{Lemma}
\newtheorem{proposition}{Proposition}
\newtheorem{corollary}{Corollary}
\newtheorem{assumption}{Assumption}
\newtheorem{remark}{Remark}

\hypersetup{
    colorlinks=true,
    linkcolor=blue!70!black,
    citecolor=green!50!black,
    urlcolor=blue!80!black
}

\title{\textbf{SIGMA-QGD: A Zero-Extra-Circuit, Variance-Aware Curvature Proxy for Variational Quantum Eigensolvers}}

\author{
    \textbf{Arjun E Naik}$^{1}$ \quad \textbf{Jitendra Tembhurne}$^{1}$ \\[1ex]
    $^{1}$Department of Computer Science, Indian Institute of Information Technology, Nagpur, India \\
    \texttt{\{bt24csd005, jtembhurne\}@iiitn.ac.in}
}
\date{}

\begin{document}

\maketitle

\begin{abstract}
Classical optimization remains the primary practical bottleneck in Variational Quantum Eigensolvers (VQEs): even when exact analytic parameter-shift gradients are available, the cost landscape exhibits severe non-stationarity, alternating between steep, ill-conditioned ravines and vast, barren-plateau-like regions. Existing optimizers either ignore curvature information entirely (e.g., Adam, SPSA) or extract it via the Quantum Fisher Information (QFI) metric at the expense of substantial additional quantum circuit evaluations per iteration (e.g., QNG requiring $3p+2$ circuits). In this work, we introduce \textbf{SIGMA-QGD}, a classical optimizer founded on a simple principle: the temporal variance of an individual parameter's gradient component across recent optimization steps serves as an effective, zero-extra-circuit surrogate for its local landscape curvature. We provide a rigorous, assumption-qualified theoretical foundation for this trajectory-variance proxy, establishing: (1) an asymptotic proportionality between trajectory variance and squared Hessian diagonal elements under local quadraticity and aggregate diagonal dominance ($\delta \sqrt{p} \ll |H_{ii}|$); (2) uniform two-sided bounds on the resulting preconditioner; (3) a descent-direction guarantee formulated under an empirical coordinate-wise descent margin condition; and (4) an ergodic $\mathcal{O}(1/T)$ non-convex stationarity convergence rate ($\mathcal{O}(1/\sqrt{T})$ under finite-shot quantum measurement noise). We validate the proxy empirically against the exact analytic parameter-shift Hessian diagonal across 10- and 12-qubit TFIM and XXZ models (Pearson $r \in [0.57, 0.63]$, $r^2 \approx 0.32\text{--}0.39$ on TFIM; $r \approx 0.30$ on XXZ; $p < 10^{-4}$ across $>394{,}000$ coordinate evaluations). Benchmarking across 11 optimizers (SIGMA-QGD alongside 10 baseline algorithms) on Transverse-Field Ising (TFIM) and Heisenberg XXZ Hamiltonians at 10 and 12 qubits demonstrates that SIGMA-QGD achieves an order-of-magnitude accuracy advantage on the commuting TFIM landscape and maintains the lowest mean energy gap on the frustrated XXZ chain while strictly adhering to a standard $2p+1$ circuit budget, statistically tying with natural gradient and adaptive baselines on 12-qubit XXZ.
\end{abstract}

\textbf{Keywords:} Variational Quantum Eigensolver (VQE), Classical Optimization, Curvature Preconditioning, Gradient Trajectory Variance, Quantum Natural Gradient.

\newpage

% ==============================================================================
% SECTION 1: INTRODUCTION
% ==============================================================================
\section{Introduction}
\label{sec:intro}

Variational Quantum Eigensolvers (VQEs) \cite{peruzzo2014variational,mcclean2016theory} represent a cornerstone of noisy intermediate-scale quantum (NISQ) computing. In a VQE, a parameterized quantum circuit (ansatz) prepares a quantum state $|\psi(\boldsymbol{\theta})\rangle$ with parameters $\boldsymbol{\theta} \in \mathbb{R}^p$. A quantum processor evaluates the expectation value of a target problem Hamiltonian $H$:
\begin{equation}
E(\boldsymbol{\theta}) = \langle \psi(\boldsymbol{\theta}) | H | \psi(\boldsymbol{\theta}) \rangle,
\label{eq:vqe_cost}
\end{equation}
while an external classical optimization algorithm iteratively updates $\boldsymbol{\theta}$ to estimate the ground-state energy $E_0 = \min_{\boldsymbol{\theta}} E(\boldsymbol{\theta})$.

\subsection{The Optimization Bottleneck in VQE}
Despite the widespread deployment of exact parameter-shift rules \cite{mitarai2018quantum,schuld2019evaluating} to obtain unbiased analytic gradients, optimizing $E(\boldsymbol{\theta})$ remains notoriously challenging. The cost landscapes of parameterized quantum circuits are characterized by severe non-stationarity \cite{mcclean2018barren,cerezo2021cost}: flat, gradient-vanishing regions (barren plateaus) coexist with narrow, highly anisotropic curved ravines \cite{mcclean2018barren,cerezo2021cost,wierichs2020avoiding}. 

Standard first-order classical optimizers fail to address these landscape characteristics adequately:
\begin{itemize}
    \item \textbf{Isotropic and Empirical Variance Optimizers (Adam, RMSProp, Yogi):} Classical adaptive gradient methods \cite{kingma2014adam,tieleman2012lecture,zaheer2018adaptive} track squared gradients or momentum-referenced variance \cite{zhuang2020adabelief}. However, they treat gradient fluctuations as isotropic and stationary across iterations, often stalling in plateaus or overshooting narrow minima.
    \item \textbf{Quantum Natural Gradient (QNG):} Metric-aware algorithms \cite{stokes2020quantum,wierichs2020avoiding} rescale gradients via the Fubini--Study Quantum Fisher Information (QFI) metric tensor. While highly effective at navigating geometry, diagonal QNG requires evaluating $(p+1)$ additional quantum circuit states per step, inflating the per-step cost from $2p+1$ to $3p+2$ circuits.
    \item \textbf{Derivative-Free Methods (SPSA, COBYLA):} Methods that bypass parameter-shift gradients \cite{spall1992multivariate,novak2025benchmarking} discard curvature information altogether, leading to slow convergence in anisotropic landscapes.
\end{itemize}

\subsection{Trajectory-Variance Curvature Preconditioning}
This paper explores whether meaningful curvature awareness can be recovered \emph{purely from the temporal statistics of standard parameter-shift gradients}, incurring \textbf{zero additional quantum circuit evaluations}. 

We propose \textbf{SIGMA-QGD} (Sampled Iterative Gradient Moment Adaptation for Quantum Gradient Descent). Rather than querying quantum hardware for quantum metric overlaps or second derivatives, SIGMA-QGD maintains an online Bessel-corrected Welford accumulator of each parameter's gradient across optimization steps. When a coordinate exhibits large temporal fluctuation along its trajectory, this instability reflects sharp local curvature, warranting step damping. Conversely, coordinates exhibiting stable, low-variance gradients are accelerated across flat regions.

\subsection{Theoretical Contributions}
The key contributions of this paper are structured as follows:
\begin{itemize}
    \item \textbf{Theoretical Foundation (Section~\ref{sec:theory}):} We formalize the relationship between temporal gradient variance $\hat{\sigma}_i^2$ and local Hessian diagonal curvature $H_{ii}^2$ (Proposition~\ref{prop:variance_curvature}), prove uniform two-sided bounds on the preconditioner (Proposition~\ref{prop:boundedness}), establish a strict descent guarantee under a coordinate-wise margin condition with dynamical self-damping (Proposition~\ref{prop:descent}), and prove $\mathcal{O}(1/T)$ noiseless and $\mathcal{O}(1/\sqrt{T})$ finite-shot non-convex stationarity convergence rates (Theorem~\ref{thm:convergence}).
    \item \textbf{Empirical Curvature Validation (Section~\ref{sec:exp_proxy}):} We validate Proposition~\ref{prop:variance_curvature} empirically by evaluating Pearson and Spearman correlations between our online variance estimator and the exact parameter-shift diagonal Hessian.
    \item \textbf{Systematic Experimental Protocol (Section~\ref{sec:setup}):} We present an experimental methodology including an independent learning-rate grid sweep, component ablations, finite-shot noise stress tests, and budget-matched comparisons against metaheuristics.
    \item \textbf{Benchmark Results (Section~\ref{sec:results}):} We present comprehensive evaluations on 10-qubit and 12-qubit TFIM and XXZ spin chains against 10 baseline optimizers (11 optimizers evaluated in total).
\end{itemize}
All implementation source code, experimental configuration files, raw per-seed trajectory datasets, and plotting pipelines are openly available on GitHub at \url{https://github.com/Arjun-E-Naik/SIGMA-QGD}.

% ==============================================================================
% SECTION 2: RELATED WORK
% ==============================================================================
\section{Related Work}
\label{sec:related}

\subsection{Adaptive First-Order Optimizers}
First-order adaptive optimizers adjust individual learning rates based on historical gradient moments. Adam \cite{kingma2014adam} normalizes the step using an exponential moving average of squared gradients $\mathbb{E}[g_i^2]$. RMSProp \cite{tieleman2012lecture} uses uncentered second moments, while Yogi \cite{zaheer2018adaptive} controls the rate of second-moment updates in non-stationary settings. AdaBelief \cite{zhuang2020adabelief} modifies Adam by computing the variance around the momentum estimate $\mathbb{E}[(g_i - m_i)^2]$. In contrast, SIGMA-QGD computes the variance around the gradient's own running mean via Welford's algorithm, avoiding conflation with the lagged momentum signal.

\subsection{Quantum Natural Gradient and Curvature Preconditioning}
Quantum Natural Gradient (QNG) \cite{stokes2020quantum} updates parameters according to $\boldsymbol{\theta}_{t+1} = \boldsymbol{\theta}_t - \eta G^{-1} \mathbf{g}_t$, where $G$ is the Fubini--Study metric tensor. While block-diagonal and full QNG require $\mathcal{O}(p^2)$ circuits, diagonal QNG requires $p+1$ metric circuits per iteration. Hamiltonian-aware QNG \cite{shi2025hamiltonian} and Stein-identity QFI estimation \cite{halla2025stein} approximate curvature using specialized circuit evaluations. 

While recent methods such as Hamiltonian-aware QNG (H-QNG)~\cite{shi2025hamiltonian} and Stein-identity QFI estimation~\cite{halla2025stein} reduce circuit overhead compared to full QNG, they differ fundamentally from SIGMA-QGD in both computational mechanics and hardware scaling. H-QNG constructs a pullback metric tensor from per-Pauli-term expectation derivatives without extra quantum circuits, but requires storing and inverting a $p \times p$ (or block-diagonal) Jacobian product across individual Hamiltonian sub-terms at each step; furthermore, it relies on ungrouped term-wise derivative access, which is highly sensitive to shot noise when non-commuting Pauli strings are sampled with finite measurements. Similarly, Stein-identity QFI estimation evaluates the metric tensor via random parameter perturbations, introducing stochastic Monte Carlo variance and step-perturbation hyperparameters similar to 2-SPSA and QN-SPSA~\cite{gacon2021simultaneous}. In contrast, SIGMA-QGD operates purely on the scalar total energy gradient vector $\mathbf{g}_t$ in $\mathcal{O}(p)$ time and memory, requires no matrix inversion or random state perturbations, and naturally regularizes finite-shot measurement variance via temporal Welford accumulation. 

For hardware-efficient ans\"atze parameterized by single-qubit Pauli rotation gates $U_i(\theta_i) = \exp(-i \frac{\theta_i}{2} P_i)$ with $P_i \in \{X, Y, Z\}$, the generator eigenvalues are $\pm 1/2$. Consequently, the energy expectation value along coordinate $\theta_i$ is a single-frequency sinusoid:
\begin{equation}
E(\theta_i) = A_i \cos(\theta_i + \phi_i) + C_i,
\label{eq:sinusoid_form}
\end{equation}
where $A_i, \phi_i, C_i$ depend on the remaining parameters $\boldsymbol{\theta}_{\ne i}$. Differentiating analytically yields gradient $g_i(\boldsymbol{\theta}) = -A_i \sin(\theta_i + \phi_i)$ and second derivative $H_{ii}(\boldsymbol{\theta}) = \frac{\partial^2 E}{\partial \theta_i^2} = -A_i \cos(\theta_i + \phi_i) = -(E(\boldsymbol{\theta}) - C_i)$. Evaluating the trigonometric identity at shifts $\pm \pi/2$ reveals that the central difference:
\begin{equation}
H_{ii}(\boldsymbol{\theta}) = \frac{1}{2}\left(E\left(\boldsymbol{\theta} + \frac{\pi}{2}\mathbf{e}_i\right) + E\left(\boldsymbol{\theta} - \frac{\pi}{2}\mathbf{e}_i\right) - 2E(\boldsymbol{\theta})\right)
\label{eq:exact_hessian}
\end{equation}
is algebraically exact \cite{wierichs2022general}. Crucially, because the forward and backward shifted evaluations $E(\boldsymbol{\theta} \pm \frac{\pi}{2}\mathbf{e}_i)$ are already executed by the quantum processor to compute the parameter-shift gradient $g_i(\boldsymbol{\theta})$ in Eq.~\eqref{eq:param_shift}, evaluating Eq.~\eqref{eq:exact_hessian} requires only the single unshifted evaluation $E(\boldsymbol{\theta})$, totaling $2p+1$ circuits. In our benchmark, we evaluate this instantaneous second-order signal directly through two independent baselines: a diagonal Hessian-magnitude baseline (\textbf{Hess-Mag}, scaling the update by $|H_{ii}| + \lambda_{\text{reg}}$) and a Levenberg--Marquardt safeguarded diagonal Newton method (\textbf{Diag-Hessian}). Crucially, Hess-Mag is evaluated purely as an instantaneous classical curvature baseline and is distinct from the quantum-Fisher-derived Hamiltonian-aware QNG of Shi \& Wang \cite{shi2025hamiltonian}.

A natural and fundamental question arises: \emph{if $H_{ii}(\boldsymbol{\theta})$ can be evaluated at $2p+1$ circuits, why estimate curvature via gradient trajectory variance $\hat{\sigma}_i^2$ rather than directly employing $H_{ii}(\boldsymbol{\theta}_t)$?} Three critical structural reasons differentiate SIGMA-QGD:
\begin{itemize}
    \item \textbf{Strict $2p$ Circuit Requirement:} Standard parameter-shift gradient estimation (Eq.~\eqref{eq:param_shift}) requires strictly $2p$ circuit evaluations. The $+1$ unshifted circuit $E(\boldsymbol{\theta})$ is required solely by Eq.~\eqref{eq:exact_hessian} to close the finite difference (or for energy logging). SIGMA-QGD operates directly on gradient vectors $\mathbf{g}_t$ and can execute in strictly $2p$ circuits per step without querying unshifted state energies.
    \item \textbf{Generality Across Multi-Qubit and Non-Pauli Generators:} Eq.~\eqref{eq:exact_hessian} yields $H_{ii}$ from single-parameter shifts only when gate generators possess two eigenvalues $\pm 1/2$. For general multi-qubit entangling gates, fermionic excitation generators (e.g., UCCSD, $k$-UpCCGSD), or QAOA mixers with multi-frequency spectra, second derivatives require $4$- or $8$-point shift rules \cite{wierichs2022general}, doubling or tripling circuit costs. In contrast, SIGMA-QGD's trajectory variance proxy remains strictly zero-extra-circuits for arbitrary variational ansatzes.
    \item \textbf{Immunity to Singular Transitions on Non-Convex Landscapes:} On trigonometric quantum landscapes, instantaneous diagonal second derivatives $H_{ii}(\boldsymbol{\theta}_t)$ oscillate rapidly, frequently switch signs, and pass through near-zero inflection points ($H_{ii} \approx 0$). In instantaneous methods (Diag-Hessian and Hess-Mag), the regularized inverse $(|H_{ii}| + \lambda_{\text{reg}})^{-1}$ spikes erratically at inflection points, provoking directional overshooting unless constrained by severe dampening. Conversely, SIGMA-QGD accumulates trajectory curvature variance across steps: $\hat{\sigma}_i^2(t)$ acts as an integrated, strictly positive-definite curvature memory that is inherently immune to point-wise Hessian singular crossings. Indeed, our ablation study (Section~\ref{sec:ablation_results}) confirms that when momentum is removed from SIGMA-QGD ($\beta=0$), performance degrades to $\Delta E \approx 0.4202$ (10q) and $0.5544$ (12q) on TFIM, directly matching Diag-Hessian ($0.4384$, $0.5609$) and Hess-Mag ($0.4431$, $0.5497$), proving that temporal trajectory integration is essential for stable preconditioning on quantum landscapes.
\end{itemize}

More recently, hybrid schemes combining exact parameter-shift gradients with simultaneous perturbation stochastic approximation (QN-SPSA+PSR)~\cite{qnspsa_psr_2024} have been explored to approximate natural gradient metrics at modest circuit costs. However, simultaneous stochastic perturbations introduce persistent Monte Carlo estimator variance that hampers asymptotic convergence on multi-parameter spin chains. Complementary investigations into directional forward-mode gradient estimators and Riemannian manifold geometry, such as QUIVER and Exact Geodesic Transport (EGT), explore alternative avenues for geometry awareness. However, directional forward gradients require randomized projection directions that introduce stochastic sampling noise, while exact geometric transport techniques depend on analytically tractable ans\"atze rather than general parameterized quantum circuits. In contrast, SIGMA-QGD operates deterministically on standard parameter-shift gradient vectors, extracting curvature purely from realized optimization trajectories.

\subsection{Quasi-Newton Methods and Line-Search Constraints}
Classical deterministic quasi-Newton algorithms (such as L-BFGS-B and non-linear Conjugate Gradient) are standard baselines in noiseless classical optimization. However, their direct deployment in physical VQE pipelines is severely constrained by quantum execution costs. Quasi-Newton secant updates and Conjugate Gradient step calculations require line searches (e.g., Armijo backtracking or Wolfe conditions) to maintain positive-definiteness ($s_k^\top y_k > 0$) and ensure descent. In VQE, each line-search trial step requires executing independent quantum state preparation and measurement circuits. An adaptive line search averaging $3\text{--}5$ trial evaluations per step inflates the per-step circuit count by $300\%\text{--}500\%$, violating fixed hardware budget parity. Furthermore, under realistic finite sampling shot noise, stochastic energy fluctuations cause Wolfe conditions to fail or cycle indefinitely. Consequently, in budget-matched VQE benchmarking under fixed evaluation constraints ($2p+1$ circuits/step), gradient-based updates with analytic step sizes or momentum schedules are standard practice.

\subsection{Positioning Summary}
Table~\ref{tab:positioning} positions SIGMA-QGD within the broader taxonomy of VQE optimizers.

\begin{table}[htbp]
\centering
\small
\caption{Positioning of representative VQE optimizers along structural and cost axes.}
\label{tab:positioning}
\resizebox{\columnwidth}{!}{%
\begin{tabular}{lcccc}
\toprule
\textbf{Method} & \textbf{Curvature-Aware} & \textbf{Extra Circuit Type} & \textbf{Circuits / Step} & \textbf{Descent Guarantee} \\
\midrule
COBYLA \cite{powell1994direct} & No & None & $1$ & No \\
SPSA \cite{spall1992multivariate} & No & None & $3$ & Stochastic \\
Adam \cite{kingma2014adam} & No & None & $2p$ (or $2p+1$) & Empirical \\
AdaBelief \cite{zhuang2020adabelief} & Partial (Mom. Var) & None & $2p$ (or $2p+1$) & Empirical \\
Diag-QNG \cite{stokes2020quantum} & Yes (QFI metric) & Metric Overlaps & $3p+2$ & Yes (Metric) \\
Hess-Mag (This work) & Yes (Exact $|H_{ii}|$) & None & $2p+1$ & Empirical \\
Diag-Hessian (This work) & Yes (Exact $H_{ii}$) & None & $2p+1$ & Yes (LM-Safeguard) \\
L-BFGS-B / CG & Yes (Secant approx.) & Line Searches & $2p + \mathcal{O}(k)$ & Line-search dependent \\
\textbf{SIGMA-QGD (Ours)} & \textbf{Yes (Trajectory Proxy)} & \textbf{None} & $\mathbf{2p\text{ (or }2p+1\text{)}}$ & \textbf{Yes (Empirical Margin)} \\
\bottomrule
\end{tabular}%
}
\end{table}

% ==============================================================================
% SECTION 3: METHODOLOGY & THEORETICAL ANALYSIS
% ==============================================================================
\section{Methodology and Theoretical Analysis}
\label{sec:theory}

\subsection{Problem Setup and Notation}
Let $H = \sum_k c_k P_k$ be a molecular or spin Hamiltonian represented as a linear combination of Pauli strings $P_k \in \{I, X, Y, Z\}^{\otimes n}$. We consider a hardware-efficient ansatz \cite{kandala2017hardware} with $n$ qubits and $r$ repeated layers. Each layer consists of parameterized single-qubit $RY(\theta_{l,j})$ rotations followed by a linear entangling chain of CNOT gates, yielding $p = n \cdot r$ trainable parameters.

Parameters are initialized i.i.d.\ from $\text{Uniform}(-\pi/8, \pi/8)$. The partial derivative with respect to parameter $\theta_i$ is evaluated via the exact parameter-shift rule:
\begin{equation}
g_i(\boldsymbol{\theta}) = \frac{1}{2} \left[ E\left(\boldsymbol{\theta} + \frac{\pi}{2}\mathbf{e}_i\right) - E\left(\boldsymbol{\theta} - \frac{\pi}{2}\mathbf{e}_i\right) \right].
\label{eq:param_shift}
\end{equation}
Throughout this paper, $\mathbf{a} \oslash \mathbf{b}$ denotes Hadamard (elementwise) division: $(\mathbf{a} \oslash \mathbf{b})_i = a_i / b_i$.

\subsection{Online Trajectory Variance Estimation}
SIGMA-QGD maintains an online Bessel-corrected Welford accumulator \cite{welford1962note} for each parameter coordinate $i \in \{1, \dots, p\}$. At step $t \ge 1$, given gradient vector $\mathbf{g}_t$:
\begin{align}
M_{t,i} &= M_{t-1,i} + \frac{g_{t,i} - M_{t-1,i}}{t}, \label{eq:welford_m} \\
S_{t,i} &= S_{t-1,i} + (g_{t,i} - M_{t-1,i})(g_{t,i} - M_{t,i}). \label{eq:welford_s}
\end{align}
For $t \ge 2$, the unbiased sample variance is given by:
\begin{equation}
\hat{\sigma}_i^2(t) = \frac{S_{t,i}}{t - 1}, \quad \text{with } \hat{\sigma}_i^2(1) \equiv 0.
\label{eq:welford_var}
\end{equation}

\subsection{Preconditioner Formulation}
The per-coordinate curvature proxy $F_i$ and preconditioned search direction $\mathbf{d}_t$ are defined as:
\begin{equation}
F_i(t) = c \cdot \hat{\sigma}_i^2(t) + \lambda_{\text{reg}}, \qquad \mathbf{d}_t = \hat{\mathbf{m}}_t \oslash \mathbf{F}(t),
\label{eq:preconditioner}
\label{eq:preconditioned_update}
\end{equation}
where $c > 0$ is the preconditioner multiplier ($c = 1.0$), $\lambda_{\text{reg}} > 0$ is the default regularization constant ($\lambda_{\text{reg}} = 0.005$, selected via sensitivity analysis to provide a robust curvature floor), and $\hat{\mathbf{m}}_t$ is the bias-corrected momentum:
\begin{equation}
\mathbf{m}_t = \beta \mathbf{m}_{t-1} + (1-\beta)\mathbf{g}_t, \quad \hat{\mathbf{m}}_t = \frac{\mathbf{m}_t}{1 - \beta^t},
\label{eq:momentum}
\end{equation}
with parameter-adaptive momentum decay:
\begin{equation}
\beta = \operatorname{clip}\Bigl(1 - \frac{1}{\sqrt{\max(p, 4)}}, \, 0.5, \, 0.98\Bigr).
\label{eq:beta_formula}
\end{equation}
As detailed in Algorithm~\ref{alg:sigma_qgd}, prior to updating Eqs.~\eqref{eq:welford_m}--\eqref{eq:welford_s} and~\eqref{eq:momentum}, $\mathbf{g}_t$ is safeguarded via global $\ell_2$-norm clipping $\mathbf{g}_{c,t} = \mathbf{g}_t \cdot \min(1, \sqrt{p}/\|\mathbf{g}_t\|_2)$, which coincides with $\mathbf{g}_t$ whenever $\|\mathbf{g}_t\|_2 \le \sqrt{p}$.

\subsection{Theoretical Guarantees}
We now formalize the conditions under which trajectory variance accurately proxies local curvature and guarantees descent.

Let $E: \mathbb{R}^p \to \mathbb{R}$ be twice continuously differentiable with gradient $\mathbf{g} = \nabla E$ and Hessian $H = \nabla^2 E$. Consider a sliding window $W_t \subset \{1, \dots, t\}$ of $w$ steps ending at iteration $t$, with temporal variance:
\begin{equation}
\hat{\sigma}_i^2(t) = \frac{1}{w-1} \sum_{\tau \in W_t} \left( g_i(\boldsymbol{\theta}_\tau) - \bar{g}_i(W_t) \right)^2, \quad \bar{g}_i(W_t) = \frac{1}{w} \sum_{\tau \in W_t} g_i(\boldsymbol{\theta}_\tau).
\end{equation}

\begin{assumption}[Local Quadraticity]
\label{ass:quadraticity}
On window $W_t$, $E$ satisfies a second-order Taylor expansion about reference point $\boldsymbol{\theta}^* \in W_t$:
\begin{equation}
\mathbf{g}(\boldsymbol{\theta}_\tau) = \mathbf{g}(\boldsymbol{\theta}^*) + H(\boldsymbol{\theta}^*)(\boldsymbol{\theta}_\tau - \boldsymbol{\theta}^*) + \mathbf{R}(\boldsymbol{\theta}_\tau), \quad \|\mathbf{R}(\boldsymbol{\theta}_\tau)\| \le \frac{L}{2} \|\boldsymbol{\theta}_\tau - \boldsymbol{\theta}^*\|^2,
\end{equation}
where $L < \infty$ is the local Lipschitz constant of the Hessian.
\end{assumption}

\begin{assumption}[Aggregate Diagonal Dominance]
\label{ass:diag_dominance}
Off-diagonal Hessian elements are bounded by $|H_{ij}(\boldsymbol{\theta}^*)| \le \delta$ for all $i \ne j$, satisfying the aggregate dominance condition $\delta \sqrt{p} \ll |H_{ii}(\boldsymbol{\theta}^*)|$ for the active coordinates.
\end{assumption}

\begin{assumption}[Bounded Gradient and Curvature]
\label{ass:bounded_curvature}
The coordinate gradient and diagonal Hessian elements are bounded: $|g_i(\boldsymbol{\theta})| \le M$ and $|H_{ii}(\boldsymbol{\theta})| \le M$ for all $\boldsymbol{\theta}$ and coordinate indices $i$. For hardware-efficient trigonometric ansatzes, $M$ is bounded by $\sum_k |c_k|$.
\end{assumption}

\begin{assumption}[Bounded Window Displacement]
\label{ass:displacement}
The trajectory displacement within window $W_t$ is bounded by $\|\boldsymbol{\theta}_\tau - \boldsymbol{\theta}^*\| \le r$ for all $\tau \in W_t$.
\end{assumption}

\begin{assumption}[Coordinate-wise Momentum Alignment]
\label{ass:alignment}
At step $t$, the bias-corrected momentum and negative gradient agree in sign coordinate-wise:
\begin{equation}
\hat{m}_{t,i} (-g_{t,i}) \ge 0 \quad \forall i \in \{1, \dots, p\}.
\end{equation}
\end{assumption}

\begin{remark}[Aggregate Alignment vs. Coordinate-wise Margin]
Aggregate alignment $\langle \hat{\mathbf{m}}_t, -\mathbf{g}_t \rangle > 0$ alone is insufficient for diagonal preconditioning. Consider a two-dimensional system with $\mathbf{F} = (0.1, 0.005)^\top$ and $\hat{\mathbf{m}}_t \odot (-\mathbf{g}_t) = (1.0, -0.5)^\top$. Aggregate alignment holds:
\[
\langle \hat{\mathbf{m}}_t, -\mathbf{g}_t \rangle = 1.0 - 0.5 = 0.5 > 0.
\]
However, the preconditioned inner product evaluates to:
\[
\langle \mathbf{d}_t, -\mathbf{g}_t \rangle = \frac{1.0}{0.1} + \frac{-0.5}{0.005} = 10 - 100 = -90 < 0.
\]
Because the coordinate with conflicting sign is divided by a small preconditioner value $F_i$, it dominates the inner product, turning the update into an ascent direction. While coordinate-wise alignment (Assumption~\ref{ass:alignment}) provides a convenient sufficient condition, the general necessary and sufficient criterion is the margin condition formalized in Proposition~\ref{prop:descent}.
\end{remark}

\begin{lemma}[Local Secant Relation]
\label{lem:secant}
Under Assumptions~\ref{ass:quadraticity}, \ref{ass:diag_dominance}, and \ref{ass:displacement}, the gradient difference satisfies:
\begin{equation}
g_i(\boldsymbol{\theta}_\tau) - g_i(\boldsymbol{\theta}^*) = H_{ii}(\boldsymbol{\theta}^*)(\theta_{\tau,i} - \theta_i^*) + \varepsilon_{i,\tau},
\end{equation}
where $|\varepsilon_{i,\tau}| \le \delta r \sqrt{p} + \frac{L}{2} r^2$.
\end{lemma}

\begin{proof}
Decomposing the Taylor expansion coordinate-wise:
\begin{equation}
g_i(\boldsymbol{\theta}_\tau) - g_i(\boldsymbol{\theta}^*) = H_{ii}(\boldsymbol{\theta}^*)(\theta_{\tau,i} - \theta_i^*) + \sum_{j \ne i} H_{ij}(\boldsymbol{\theta}^*)(\theta_{\tau,j} - \theta_j^*) + R_i(\boldsymbol{\theta}_\tau).
\end{equation}
By Cauchy--Schwarz and Assumption~\ref{ass:diag_dominance}:
\begin{equation}
\left| \sum_{j \ne i} H_{ij}(\boldsymbol{\theta}^*)(\theta_{\tau,j} - \theta_j^*) \right| \le \delta \sum_{j=1}^p |\theta_{\tau,j} - \theta_j^*| \le \delta \sqrt{p} \|\boldsymbol{\theta}_\tau - \boldsymbol{\theta}^*\| \le \delta r \sqrt{p}.
\end{equation}
Combining this with $\|R_i(\boldsymbol{\theta}_\tau)\| \le \|\mathbf{R}(\boldsymbol{\theta}_\tau)\| \le \frac{L}{2} r^2$ yields the stated bound.
\end{proof}

\begin{proposition}[Variance--Curvature Proportionality]
\label{prop:variance_curvature}
Under Assumptions~\ref{ass:quadraticity}--\ref{ass:displacement}, the sample variance of the $i$-th gradient component on window $W_t$ satisfies:
\begin{equation}
\hat{\sigma}_i^2(t) = H_{ii}(\boldsymbol{\theta}^*)^2 \operatorname{Var}_{W_t}[\theta_i] + \mathcal{R}_i,
\end{equation}
where the residual $\mathcal{R}_i$ is rigorously bounded by:
\begin{equation}
|\mathcal{R}_i| \le 2 M r \left( \delta r \sqrt{p} + \frac{L}{2} r^2 \right) + \left( \delta r \sqrt{p} + \frac{L}{2} r^2 \right)^2 = \mathcal{O}\left( (M \delta \sqrt{p} + \delta^2 p) r^2 + M L r^3 + L^2 r^4 \right).
\label{eq:prop1_remainder}
\end{equation}
Consequently, when the parameter moves actively across the window ($\operatorname{Var}_{W_t}[\theta_i] = \Theta(r^2)$), the diagonal curvature term $H_{ii}(\boldsymbol{\theta}^*)^2 \operatorname{Var}_{W_t}[\theta_i]$ dominates the residual provided $\delta \sqrt{p} \ll |H_{ii}(\boldsymbol{\theta}^*)|$ (Assumption~\ref{ass:diag_dominance}) and displacement $r$ is sufficiently small.
\end{proposition}

\begin{proof}
Applying the variance operator to the secant relation from Lemma~\ref{lem:secant}:
\begin{equation}
\operatorname{Var}_{W_t}[g_i] = \operatorname{Var}_{W_t}\left[ H_{ii}(\boldsymbol{\theta}^*) \theta_i + \varepsilon_i \right] = H_{ii}(\boldsymbol{\theta}^*)^2 \operatorname{Var}_{W_t}[\theta_i] + 2 H_{ii}(\boldsymbol{\theta}^*) \operatorname{Cov}_{W_t}(\theta_i, \varepsilon_i) + \operatorname{Var}_{W_t}[\varepsilon_i].
\end{equation}
By the Cauchy--Schwarz inequality, $|\operatorname{Cov}_{W_t}(\theta_i, \varepsilon_i)| \le \sqrt{\operatorname{Var}_{W_t}[\theta_i]} \sqrt{\operatorname{Var}_{W_t}[\varepsilon_i]}$. Because the displacement satisfies $\|\boldsymbol{\theta}_\tau - \boldsymbol{\theta}^*\| \le r$, we have $|\theta_{\tau,i} - \theta_i^*| \le r$, which implies $\sqrt{\operatorname{Var}_{W_t}[\theta_i]} \le r$ (up to the finite-window Bessel factor $\sqrt{\frac{w}{w-1}} = 1 + \mathcal{O}(w^{-1})$). From Lemma~\ref{lem:secant}, $\|\boldsymbol{\varepsilon}_i\|_\infty \le \delta r \sqrt{p} + \frac{L}{2} r^2$, yielding $\sqrt{\operatorname{Var}_{W_t}[\varepsilon_i]} \le \|\boldsymbol{\varepsilon}_i\|_\infty$ (up to $1 + \mathcal{O}(w^{-1})$). Using $|H_{ii}(\boldsymbol{\theta}^*)| \le M$, the cross-covariance term is bounded by:
\begin{equation}
2 |H_{ii}(\boldsymbol{\theta}^*)| |\operatorname{Cov}_{W_t}(\theta_i, \varepsilon_i)| \le 2 M r \left( \delta r \sqrt{p} + \frac{L}{2} r^2 \right).
\end{equation}
Adding $\operatorname{Var}_{W_t}[\varepsilon_i] \le \|\boldsymbol{\varepsilon}_i\|_\infty^2 = (\delta r \sqrt{p} + \frac{L}{2} r^2)^2$ completes the proof of Eq.~\eqref{eq:prop1_remainder}.
\end{proof}

\begin{remark}[Dynamical Variance Scaling, Energy Rescaling, and Window Formulation]
\label{rem:variance_scaling}
Three theoretical nuances warrant explicit clarification:
\begin{itemize}
    \item \textbf{Curvature vs. Trajectory Dampening and Hamiltonian Energy Scale:} In classical Newton-Raphson methods, steps scale inversely with the first power of curvature ($H_{ii}^{-1}$). In SIGMA-QGD, the variance proxy links to $H_{ii}^2 \operatorname{Var}[\theta_i]$, meaning steps scale approximately as $\hat{m}_{t,i} / (c H_{ii}^2 \operatorname{Var}[\theta_i] + \lambda_{\text{reg}})$. Along coordinates where steep curvature induces rapid oscillations, the denominator scales quadratically, quenching unstable dynamics. Conversely, along flat coordinates ($H_{ii} \approx 0$), the step is safely bounded by $\lambda_{\text{reg}}^{-1}$. Notice that because $d_i \sim \hat{m}_i / \hat{\sigma}_i^2 \propto 1/g_i$ in the variance-dominated regime, the update direction is not invariant under arbitrary scalar rescaling of the Hamiltonian ($H \to \alpha H$). This reflects the natural physics of unnormalized second-order updates, tying optimal choices of $c$ and $\lambda_{\text{reg}}$ to the characteristic energy scale $\sum_k |c_k|$ of the target Hamiltonian. For general Hamiltonians with arbitrary spectral norm $\Lambda = \sum_k |c_k|$, scale invariance is effortlessly restored either by normalizing the cost function $\tilde{E}(\boldsymbol{\theta}) = E(\boldsymbol{\theta}) / \Lambda$ prior to optimization or by scaling the regularization floor proportionally as $\lambda_{\text{reg}} \propto \Lambda^2$.
    \item \textbf{Trajectory-Dependent Displacement:} The main term $H_{ii}^2 \operatorname{Var}_{W_t}[\theta_i]$ depends on the coordinate variance $\operatorname{Var}_{W_t}[\theta_i]$ traversed by the optimizer itself. Thus, $\hat{\sigma}_i^2$ reflects local curvature convolved with the dynamical step history rather than an isolated, static landscape tensor.
    \item \textbf{Sliding Window vs. Online Accumulating Welford Estimator:} The theoretical derivations in Assumptions~\ref{ass:quadraticity}--\ref{ass:displacement} and Proposition~\ref{prop:variance_curvature} are formulated over a compact sliding window $W_t$ of width $w$, representing the active traversal phase ($t \in [1, t_0]$) during which the optimizer descends through an energy ravine with significant coordinate displacement ($r = \Theta(1)$). In practice, Algorithm~\ref{alg:sigma_qgd} executes an online Bessel-corrected Welford accumulator over the full historical run ($t \ge 2$) for computational and memory parsimony ($\mathcal{O}(1)$ time and space per step). As optimization approaches an asymptotic local minimum basin ($t > t_0$), parameter displacement naturally contracts ($\|\boldsymbol{\theta}_\tau - \boldsymbol{\theta}^*\| \to 0$). In a pure sliding window of width $w$, the parameter displacement variance $\operatorname{Var}_{W_t}[\theta_i]$ would decay to zero, driving $\hat{\sigma}_i^2 \to 0$ and collapsing the preconditioner to the scalar regularization floor $\lambda_{\text{reg}}$. In contrast, the full-history Welford sum-of-squares accumulator $\mathbf{S}_t / (t-1)$ retains the cumulative temporal statistics of the principal curvature modes traversed during the descent. Thus, full-history Welford prevents late-stage preconditioner collapse, maintaining directional curvature conditioning even as instantaneous parameter velocities diminish.
\end{itemize}
\end{remark}

\begin{proposition}[Uniform Local and Global Boundedness of the Preconditioner]
\label{prop:boundedness}
Under Assumptions~\ref{ass:bounded_curvature} and \ref{ass:displacement}, the diagonal preconditioner $F_i(t) = c \hat{\sigma}_i^2(t) + \lambda_{\text{reg}}$ on window $W_t$ satisfies:
\begin{equation}
\lambda_{\text{reg}} \le F_i(t) \le c M^2 r^2 + \lambda_{\text{reg}} + \mathcal{O}(r^2).
\end{equation}
Furthermore, for the full-history Welford accumulator in Algorithm~\ref{alg:sigma_qgd}, gradient clipping ($\|\mathbf{g}_{c,t}\|_2 \le c_{\text{norm}} = \sqrt{p}$) and Assumption~\ref{ass:bounded_curvature} ensure $|g_{c,t,i}| \le G_{\max} \equiv \min(\sqrt{p}, M)$ for all $t \ge 1$, yielding the unconditional global bound $\lambda_{\text{reg}} \le F_i(t) \le F_{\max} \equiv 2c G_{\max}^2 + \lambda_{\text{reg}}$ for all $t \ge 1$, and coordinate-wise step bounds:
\begin{equation}
\frac{|\hat{m}_{t,i}|}{F_{\max}} \le |d_{t,i}| \le \frac{|\hat{m}_{t,i}|}{\lambda_{\text{reg}}}.
\end{equation}
\end{proposition}

\begin{proof}
Since $\hat{\sigma}_i^2 \ge 0$, the lower bound follows from $\lambda_{\text{reg}} > 0$. The local window upper bound follows from Proposition~\ref{prop:variance_curvature} and $\operatorname{Var}_{W_t}[\theta_i] \le \frac{w}{w-1} r^2$ (absorbing the Bessel factor $1 + \mathcal{O}(w^{-1})$). Globally, since $|g_{c,t,i}| \le G_{\max}$, the unbiased sample variance over any horizon $t \ge 2$ satisfies $\hat{\sigma}_i^2(t) = \frac{1}{t-1}\sum_{\tau=1}^t (g_{c,\tau,i} - M_{t,i})^2 \le \frac{t}{t-1} G_{\max}^2 \le 2 G_{\max}^2$, establishing the unconditional upper bound $F_{\max}$.
\end{proof}

\begin{proposition}[Descent Direction Guarantee under Margin Condition]
\label{prop:descent}
Let $\mathcal{S}_+ = \{i : \hat{m}_{t,i}(-g_{t,i}) \ge 0\}$ and $\mathcal{S}_- = \{j : \hat{m}_{t,j}(-g_{t,j}) < 0\}$ partition the coordinate indices $\{1, \dots, p\}$ at step $t$. Under Assumptions~\ref{ass:bounded_curvature} and \ref{ass:displacement}, the preconditioned update direction $\mathbf{d}_t = \hat{\mathbf{m}}_t \oslash \mathbf{F}(t)$ constitutes a strict descent direction ($\langle \mathbf{d}_t, -\mathbf{g}_t \rangle > 0$) if the aligned-coordinate projection dominates the misaligned contribution according to the margin condition:
\begin{equation}
\sum_{i \in \mathcal{S}_+} \frac{\hat{m}_{t,i} (-g_{t,i})}{F_i(t)} > \sum_{j \in \mathcal{S}_-} \frac{|\hat{m}_{t,j} g_{t,j}|}{F_j(t)}.
\label{eq:margin_condition}
\end{equation}
In particular, using the uniform bounds $\lambda_{\text{reg}} \le F_k(t) \le F_{\max}$ from Proposition~\ref{prop:boundedness}, a sufficient condition is:
\begin{equation}
\sum_{i \in \mathcal{S}_+} \hat{m}_{t,i} (-g_{t,i}) > \kappa \sum_{j \in \mathcal{S}_-} |\hat{m}_{t,j} g_{t,j}|, \qquad \kappa = \frac{F_{\max}}{\lambda_{\text{reg}}}.
\label{eq:kappa_margin}
\end{equation}
In the special case of coordinate-wise alignment ($\mathcal{S}_- = \emptyset$, Assumption~\ref{ass:alignment}), descent holds unconditionally with $\langle \mathbf{d}_t, -\mathbf{g}_t \rangle \ge \frac{1}{F_{\max}} \langle \hat{\mathbf{m}}_t, -\mathbf{g}_t \rangle > 0$.
\end{proposition}

\begin{proof}
Decomposing the preconditioned inner product over the partition $\mathcal{S}_+ \cup \mathcal{S}_-$:
\begin{align}
\langle \mathbf{d}_t, -\mathbf{g}_t \rangle &= \sum_{i=1}^p \frac{\hat{m}_{t,i}(-g_{t,i})}{F_i(t)} \nonumber \\
&= \sum_{i \in \mathcal{S}_+} \frac{\hat{m}_{t,i}(-g_{t,i})}{F_i(t)} - \sum_{j \in \mathcal{S}_-} \frac{|\hat{m}_{t,j} g_{t,j}|}{F_j(t)}.
\end{align}
Positivity is guaranteed whenever the first sum exceeds the second, establishing Eq.~\eqref{eq:margin_condition}. Applying the uniform lower bound $\lambda_{\text{reg}}$ and upper bound $F_{\max}$ yields Eq.~\eqref{eq:kappa_margin}. When $\mathcal{S}_- = \emptyset$, the negative sum vanishes and the uniform lower bound preserves strict descent.
\end{proof}

\paragraph{Empirical Validity and Dynamical Self-Damping.}
While ideal coordinate-wise alignment ($\mathcal{S}_- = \emptyset$) holds in $67.10\%$ (TFIM) and $71.93\%$ (XXZ) of coordinate-step pairs, the generalized margin condition in Eq.~\eqref{eq:margin_condition} reveals an intrinsic self-stabilizing mechanism of trajectory-variance preconditioning. Transient sign misalignments ($\hat{m}_{t,j} g_{t,j} > 0$) occur almost exclusively along coordinates undergoing rapid directional oscillations across steep ravine walls. Along these jittery coordinates, the temporal variance $\hat{\sigma}_j^2$ is large, automatically inflating the preconditioner denominator $F_j(t) \gg \lambda_{\text{reg}}$. Consequently, the misaligned negative terms in Eq.~\eqref{eq:margin_condition} are dynamically suppressed by their large preconditioner values ($1/F_j(t) \to 0$). Conversely, coordinates progressing smoothly along the valley floor exhibit stable gradients, keeping $F_i(t)$ near $\lambda_{\text{reg}}$ and amplifying positive descent progress. Furthermore, along realized trajectories across all tested models, the coordinate-wise condition ratio $\max_{i,t} F_i(t) / \lambda_{\text{reg}}$ remains empirically bounded in $[12.4, 48.6]$, far below the theoretical worst-case ceiling $\kappa = F_{\max}/\lambda_{\text{reg}} \le 200\text{--}400$, as coordinates rarely reach theoretical clipping limits simultaneously. As a direct consequence of this dynamical self-damping, the aggregate preconditioned inner product $\langle \mathbf{d}_t, -\mathbf{g}_t \rangle$ remained strictly positive in \textbf{100.0\%} of evaluated optimization steps ($200/200$) across all tested models and system sizes.

\begin{theorem}[Non-Convex Stationarity Convergence Rate]
\label{thm:convergence}
Let $E: \mathbb{R}^p \to \mathbb{R}$ be $L_g$-smooth with lower-bounded global minimum $E^* > -\infty$. Suppose that along the optimization trajectory, the margin condition of Proposition~\ref{prop:descent} holds with uniform alignment factor $\gamma_0 > 0$ such that $\langle \mathbf{d}_t, -\mathbf{g}_t \rangle \ge \frac{\gamma_0}{F_{\max}} \|\mathbf{g}_t\|_2^2$ (which holds unconditionally when $\beta = 0$ with $\gamma_0 = \min(1, 1/M)$, or empirically under Proposition~\ref{prop:descent} for $\beta > 0$), where $F_{\max}$ is the uniform preconditioner upper bound from Proposition~\ref{prop:boundedness}, and the base learning rate satisfies $\eta_0 < \frac{2 \gamma_0 \lambda_{\text{reg}}^2 (1-\beta)}{L_g F_{\max}}$. Then the sequence of iterates $\{\boldsymbol{\theta}_{t-1}\}_{t=1}^T$ generated by Algorithm~\ref{alg:sigma_qgd} satisfies the ergodic stationary-point convergence bound:
\begin{equation}
\min_{1 \le t \le T} \|\nabla E(\boldsymbol{\theta}_{t-1})\|_2^2 \le \frac{E(\boldsymbol{\theta}_0) - E^*}{\mathcal{C} \cdot T} = \mathcal{O}\left(\frac{1}{T}\right),
\label{eq:ergodic_rate}
\end{equation}
where $\mathcal{C} = \eta_0 \left( \frac{\gamma_0}{F_{\max}} - \frac{L_g \eta_0}{2 \lambda_{\text{reg}}^2 (1-\beta)} \right) > 0$. Furthermore, in the presence of unbiased stochastic quantum measurement shot noise with bounded variance $\mathbb{E}[\|\tilde{\mathbf{g}}_t - \nabla E(\boldsymbol{\theta}_{t-1})\|_2^2 \mid \mathcal{F}_{t-1}] \le \sigma_{\text{shot}}^2$, setting a diminishing step-size schedule $\eta_t = \eta_0 / \sqrt{t}$ yields the standard non-convex stochastic stationarity rate:
\begin{equation}
\min_{1 \le t \le T} \mathbb{E}\left[\|\nabla E(\boldsymbol{\theta}_{t-1})\|_2^2\right] = \mathcal{O}\left(\frac{1}{\sqrt{T}}\right).
\label{eq:stochastic_rate}
\end{equation}
\end{theorem}

\begin{proof}
By $L_g$-smoothness of $E(\boldsymbol{\theta})$, for consecutive parameter iterates $\boldsymbol{\theta}_t = \boldsymbol{\theta}_{t-1} - \eta_0 \mathbf{d}_t$ with $\mathbf{g}_t = \nabla E(\boldsymbol{\theta}_{t-1})$:
\begin{equation}
E(\boldsymbol{\theta}_t) \le E(\boldsymbol{\theta}_{t-1}) - \eta_0 \langle \mathbf{g}_t, \mathbf{d}_t \rangle + \frac{L_g \eta_0^2}{2} \|\mathbf{d}_t\|_2^2.
\end{equation}
Applying the margin lower bound $\langle \mathbf{d}_t, -\mathbf{g}_t \rangle \ge \frac{\gamma_0}{F_{\max}} \|\mathbf{g}_t\|_2^2$ and the preconditioner lower bound $F_i(t) \ge \lambda_{\text{reg}}$ (for $t \ge t_{\text{warm}}$, and $F_i(t) = 1 \ge \lambda_{\text{reg}}$ during warmup) gives $\|\mathbf{d}_t\|_2 \le \|\hat{\mathbf{m}}_t\|_2 / \lambda_{\text{reg}}$. Summing over $t = 1, \dots, T$ yields:
\begin{equation}
E(\boldsymbol{\theta}_T) - E(\boldsymbol{\theta}_0) \le -\frac{\eta_0 \gamma_0}{F_{\max}} \sum_{t=1}^T \|\mathbf{g}_t\|_2^2 + \frac{L_g \eta_0^2}{2 \lambda_{\text{reg}}^2} \sum_{t=1}^T \|\hat{\mathbf{m}}_t\|_2^2.
\end{equation}
Since $\hat{\mathbf{m}}_t = \sum_{\tau=1}^t w_{t,\tau} \mathbf{g}_{c,\tau}$ with convex weights $w_{t,\tau} = \frac{(1-\beta)\beta^{t-\tau}}{1-\beta^t}$ summing to $1$, $\|\mathbf{g}_{c,\tau}\|_2 \le \|\mathbf{g}_\tau\|_2$, and $1 - \beta^t \ge 1 - \beta$, Jensen's inequality gives $\|\hat{\mathbf{m}}_t\|_2^2 \le \sum_{\tau=1}^t w_{t,\tau} \|\mathbf{g}_\tau\|_2^2 \le \sum_{\tau=1}^t \beta^{t-\tau} \|\mathbf{g}_\tau\|_2^2$. Summing over $t=1,\dots,T$ and exchanging summation order yields $\sum_{t=1}^T \|\hat{\mathbf{m}}_t\|_2^2 \le \frac{1}{1-\beta} \sum_{t=1}^T \|\mathbf{g}_t\|_2^2$. Substituting this bound and rearranging with $E(\boldsymbol{\theta}_T) \ge E^*$ gives:
\begin{equation}
\mathcal{C} \sum_{t=1}^T \|\mathbf{g}_t\|_2^2 \le E(\boldsymbol{\theta}_0) - E^*,
\end{equation}
where $\mathcal{C} > 0$ under the stated step-size condition. Bounding $\sum_{t=1}^T \|\mathbf{g}_t\|_2^2 \ge T \min_{1 \le t \le T} \|\mathbf{g}_t\|_2^2$ and dividing by $\mathcal{C} T$ establishes Eq.~\eqref{eq:ergodic_rate}. Under additive zero-mean shot noise with bounded variance $\sigma_{\text{shot}}^2$, note that the Welford update satisfies $|\hat{\sigma}_i^2(t) - \hat{\sigma}_i^2(t-1)| = \mathcal{O}(G_{\max}^2 / t)$, so replacing $F_i(t)$ with the $\mathcal{F}_{t-1}$-measurable lagged preconditioner $F_i(t-1)$ incurs at most an $\mathcal{O}(\eta_t / t)$ bias per step (whose sum $\sum_{t=1}^T \eta_t / t < \infty$ is bounded), while conditional expectations introduce an $\mathcal{O}(L_g \eta_t^2 \sigma_{\text{shot}}^2)$ variance penalty per step; applying $\eta_t = \eta_0 / \sqrt{t}$ yields $\min_{1 \le t \le T} \mathbb{E}[\|\nabla E(\boldsymbol{\theta}_{t-1})\|_2^2] = \mathcal{O}(1/\sqrt{T})$~\cite{ghadimi2013stochastic}.
\end{proof}

\subsection{Warmup Schedule and Complete Algorithm}
To prevent noisy variance estimates during early steps from destabilizing optimization, preconditioning is activated only after a warmup period:
\begin{equation}
t_{\text{warm}} = \min(3p, \max(10, \lfloor T / 5 \rfloor)).
\label{eq:warmup}
\end{equation}
For $t < t_{\text{warm}}$, the update is standard momentum-accelerated gradient descent ($\mathbf{d}_t = \hat{\mathbf{m}}_t$). (Evaluating the unshifted energy $E_t = E(\boldsymbol{\theta}_{t-1})$ adds 1 circuit per step for best-iterate tracking, bringing the logged experimental budget to $2p+1$; if only parameter updates are required, $E_t$ can be omitted to execute in strictly $2p$ circuits per step).

\begin{algorithm}[H]
\caption{SIGMA-QGD Optimization Algorithm}
\label{alg:sigma_qgd}
\begin{algorithmic}[1]
\REQUIRE Initial parameters $\boldsymbol{\theta}_0 \in \mathbb{R}^p$, learning rate $\eta_0$, multiplier $c=1.0$, regularizer $\lambda_{\text{reg}}=0.005$, max steps $T$.
\STATE Initialize accumulators $\mathbf{M} \leftarrow \mathbf{0}, \mathbf{S} \leftarrow \mathbf{0}, \mathbf{m} \leftarrow \mathbf{0}$, best energy $E^* \leftarrow \infty, \boldsymbol{\theta}^* \leftarrow \boldsymbol{\theta}_0$.
\STATE Compute $t_{\text{warm}} = \min(3p, \max(10, \lfloor T/5 \rfloor))$, clip threshold $c_{\text{norm}} = \sqrt{p}$.
\FOR{$t = 1$ \TO $T$}
    \STATE Evaluate cost and parameter-shift gradient: $E_t, \mathbf{g}_t \leftarrow \text{CostAndGradient}(\boldsymbol{\theta}_{t-1})$ \hfill \COMMENT{$2p+1$ circuits}
    \STATE Gradient clipping: $\mathbf{g}_c \leftarrow \mathbf{g}_t \cdot \min(1, c_{\text{norm}} / \|\mathbf{g}_t\|_2)$
    \STATE Update Welford accumulators $(\mathbf{M}, \mathbf{S})$ with $\mathbf{g}_c$ via Eqs.~\eqref{eq:welford_m}--\eqref{eq:welford_s}
    \STATE Compute unbiased variance $\hat{\boldsymbol{\sigma}}^2 \leftarrow \mathbf{S} / (t-1)$ for $t \ge 2$, else $\mathbf{0}$
    \STATE Update momentum: $\mathbf{m} \leftarrow \beta \mathbf{m} + (1-\beta)\mathbf{g}_c$, $\hat{\mathbf{m}} \leftarrow \mathbf{m} / (1-\beta^t)$
    \IF{$t \ge t_{\text{warm}}$}
        \STATE $\mathbf{F} \leftarrow c \hat{\boldsymbol{\sigma}}^2 + \lambda_{\text{reg}} \mathbf{1}$
        \STATE $\mathbf{d}_t \leftarrow \hat{\mathbf{m}} \oslash \mathbf{F}$
    \ELSE
        \STATE $\mathbf{d}_t \leftarrow \hat{\mathbf{m}}$
    \ENDIF
    \STATE $\boldsymbol{\theta}_t \leftarrow \boldsymbol{\theta}_{t-1} - \eta_0 \mathbf{d}_t$
    \IF{$E_t < E^*$}
        \STATE $E^* \leftarrow E_t, \boldsymbol{\theta}^* \leftarrow \boldsymbol{\theta}_{t-1}$
    \ENDIF
\ENDFOR
\RETURN Optimal parameters $\boldsymbol{\theta}^*$, minimum energy $E^*$
\end{algorithmic}
\end{algorithm}

% ==============================================================================
% SECTION 4: EXPERIMENTAL SETUP
% ==============================================================================
\section{Experimental Methodology}
\label{sec:setup}

\subsection{Benchmark Hamiltonians}
We evaluate optimizers on two representative quantum spin models under periodic boundary conditions:
\begin{itemize}
    \item \textbf{1D Transverse-Field Ising Model (TFIM):}
    \begin{equation}
    H_{\text{TFIM}} = -J \sum_{i=1}^n Z_i Z_{i+1} - h \sum_{i=1}^n X_i, \quad J = 1.0, \, h = 0.5.
    \end{equation}
    Exact ground states: $E_0 = -10.6356$ ($n=10$), $E_0 = -12.7626$ ($n=12$).
    \item \textbf{1D Heisenberg XXZ Spin Chain:}
    \begin{equation}
    H_{\text{XXZ}} = J_{xy} \sum_{i=1}^n (X_i X_{i+1} + Y_i Y_{i+1}) + J_z \sum_{i=1}^n Z_i Z_{i+1} + h \sum_{i=1}^n Z_i,
    \end{equation}
    with $J_{xy} = 1.0, J_z = 0.5, h = 0.1$. Exact ground states: $E_0 = -15.2761$ ($n=10$), $E_0 = -18.2291$ ($n=12$).
\end{itemize}

\subsection{Ansatz and Initialization}
We employ a hardware-efficient ansatz with depth $r=3$, producing $p=30$ trainable parameters at $n=10$ and $p=36$ at $n=12$. For each test configuration, parameter vectors $\boldsymbol{\theta}_0$ are sampled deterministically across identical random seeds ($\theta_{0,i} \sim \text{Uniform}(-\pi/8, \pi/8)$), ensuring paired comparisons across all algorithms.

\subsection{Baseline Algorithms and Circuit Accounting}
All 11 optimizers (SIGMA-QGD alongside 10 baselines) are evaluated under strict circuit accounting rules:
\begin{itemize}
    \item \textbf{First-order baselines ($2p+1$ circuits/step):} Adam \cite{kingma2014adam}, AdaBelief \cite{zhuang2020adabelief}, RMSProp \cite{tieleman2012lecture}, and Yogi \cite{zaheer2018adaptive}.
    \item \textbf{Exact second-order baselines ($2p+1$ circuits/step):} Hess-Mag and Diag-Hessian (via Eq.~\eqref{eq:exact_hessian} \cite{wierichs2022general}).
    \item \textbf{Quantum Natural Gradient ($3p+2$ circuits/step):} Diag-QNG \cite{stokes2020quantum,wierichs2020avoiding}.
    \item \textbf{Derivative-free baselines:} COBYLA ($1$ call/step \cite{powell1994direct}), SPSA ($3$ calls/step \cite{spall1992multivariate,spall1997accelerated,spall1998implementation,spall2000adaptive}), QN-SPSA ($9$ calls/step \cite{gacon2021simultaneous}), and CMA-ES (budget-matched \cite{novak2025benchmarking}).
\end{itemize}

\subsection{Statistical Significance Testing}
All reported performance differences between SIGMA-QGD and baseline methods are evaluated via two-sided paired Wilcoxon signed-rank tests. To protect against False Discovery Rate inflation across the 10 simultaneous comparisons, we apply Holm--Bonferroni stepdown corrections to all reported $p$-values.

% ==============================================================================
% SECTION 5: EMPIRICAL RESULTS
% ==============================================================================
\section{Empirical Results}
\label{sec:results}

\subsection{Curvature Proxy Validation}
\label{sec:exp_proxy}
To empirically evaluate Proposition~\ref{prop:variance_curvature}, we measure the step-by-step correlation between the log trajectory variance $\log \hat{\sigma}_i^2$ and the log exact analytic diagonal Hessian magnitude $\log |H_{ii}|$ along 200-step SIGMA-QGD trajectories across both TFIM and XXZ Hamiltonians at 10 and 12 qubits ($N=15$ seeds, yielding $>394{,}000$ coordinate-step evaluations). At each optimization step, the exact diagonal Hessian is evaluated via the trigonometric parameter-shift identity in Eq.~\eqref{eq:exact_hessian} without perturbing the optimizer trajectory.

\begin{table*}[htbp]
\centering
\small
\caption{Consolidated curvature-proxy validation across physical Hamiltonians and system sizes ($N=15$ independent seeds, 200 steps). Reported values indicate mean $\pm$ standard deviation for log-space Pearson correlation $r$ and Spearman rank correlation $\rho$, accompanied by Wilcoxon signed-rank tests against $H_0: \text{correlation} \le 0$.}
\label{tab:hessian_corr}
\resizebox{\textwidth}{!}{%
\begin{tabular}{lcccccc}
\toprule
\textbf{Model} & \textbf{System Size} & \textbf{Evaluations/Seed} & \textbf{Pearson $r$ (Log-Space)} & \textbf{Wilcoxon $p$} & \textbf{Spearman $\rho$} & \textbf{Wilcoxon $p$} \\
\midrule
TFIM (Longitudinal/Transverse) & 10 qubits ($p=30$) & 5,970 & 0.5665 $\pm$ 0.0577 & $6.1 \times 10^{-5}$ & 0.3715 $\pm$ 0.0998 & $6.1 \times 10^{-5}$ \\
TFIM (Longitudinal/Transverse) & 12 qubits ($p=36$) & 7,164 & 0.6256 $\pm$ 0.0301 & $6.1 \times 10^{-5}$ & 0.4474 $\pm$ 0.0745 & $6.1 \times 10^{-5}$ \\
XXZ (Heisenberg Exchange)      & 10 qubits ($p=30$) & 5,970 & 0.3196 $\pm$ 0.0829 & $6.1 \times 10^{-5}$ & 0.1162 $\pm$ 0.1416 & $0.010$ \\
XXZ (Heisenberg Exchange)      & 12 qubits ($p=36$) & 7,164 & 0.2938 $\pm$ 0.0914 & $6.1 \times 10^{-5}$ & 0.1113 $\pm$ 0.1404 & $0.004$ \\
\bottomrule
\end{tabular}%
}
\end{table*}

\begin{figure}[htbp]
\centering
\includegraphics[width=0.92\linewidth]{figures/fig_hessian_correlation_consolidated.pdf}
\caption{Consolidated curvature-proxy validation across 15 seeds on 10-qubit ($p=30$) and 12-qubit ($p=36$) TFIM and XXZ spin models. (a) Log-space Pearson correlation $r$ between gradient trajectory variance $\log \hat{\sigma}_i^2$ and exact diagonal Hessian $\log |H_{ii}|$. (b) Spearman rank correlation $\rho$. Diamond markers denote sample means, horizontal bars indicate medians, and individual points depict random seeds. All correlations are statistically positive ($p < 10^{-4}$ for Pearson across all models).}
\label{fig:hessian_corr_consolidated}
\end{figure}

The empirical findings presented in Table~\ref{tab:hessian_corr} and Figure~\ref{fig:hessian_corr_consolidated} demonstrate the following physical characteristics:
\begin{itemize}
    \item \textbf{High-Fidelity Tracking in Commuting-Field Regimes (TFIM):} On the TFIM spin chain, the trajectory variance exhibits a strong log-linear correlation with the exact diagonal Hessian ($r = 0.5665 \pm 0.0577$ at 10 qubits, increasing to $r = 0.6256 \pm 0.0301$ at 12 qubits, with both $p = 6.1 \times 10^{-5}$). The increase in correlation with system size demonstrates that the variance estimator scales favourably, capturing local curvature even as the Hilbert space dimension quadruples ($2^{10} \to 2^{12}$).
    \item \textbf{Statistically Significant Signal under Frustrated Exchange (XXZ):} On the non-commuting XXZ Heisenberg model, Pearson correlation remains strictly positive and statistically significant ($r = 0.3196 \pm 0.0829$ at 10 qubits, $r = 0.2938 \pm 0.0914$ at 12 qubits, $p = 6.1 \times 10^{-5}$). Spearman rank correlation is likewise positive ($\rho = 0.1162$, $p = 0.010$ at 10 qubits; $\rho = 0.1113$, $p = 0.004$ at 12 qubits).
    \item \textbf{Physical Origin of Model Discrepancy:} The disparity in correlation magnitude between TFIM ($r \approx 0.60$) and XXZ ($r \approx 0.31$) directly mirrors the underlying Hamiltonian operator structure. In TFIM, nearest-neighbour $Z_i Z_{i+1}$ interactions commute with each other and are dominated along the coordinate axes by the transverse field $X_i$, keeping off-diagonal Hessian couplings modest ($\delta \sqrt{p} \ll |H_{ii}|$, aligning with Assumption~\ref{ass:diag_dominance}). In XXZ, the non-commuting exchange terms $X_i X_{i+1} + Y_i Y_{i+1}$ induce strong off-diagonal parameter cross-talk that bounds the fraction of energy landscape curvature accessible via purely diagonal coordinate projections. This aligns with analytical studies of Hessian entry scaling in variational quantum algorithms~\cite{hessian_resolvability_2026}, which establish that non-commuting Hamiltonian terms suppress resolvable diagonal curvature modes and amplify off-diagonal couplings across the parameter manifold.
    \item \textbf{Scale Invariance of Log-Space Correlation:} While the preconditioned step in Eq.~\eqref{eq:preconditioned_update} uses $\hat{\sigma}_i^2$ directly and Proposition~\ref{prop:variance_curvature} connects $\hat{\sigma}_i^2 \propto H_{ii}^2$, proxy validation is evaluated in log-space ($\log \hat{\sigma}_i^2$ vs. $\log |H_{ii}|$). In periodic quantum circuits, coordinate curvatures and gradient variances naturally span multiple orders of magnitude (from $10^{-4}$ near plateaus to $10^0$ near steep ravines). Linear Pearson correlations on such heavy-tailed distributions are heavily distorted by extreme outlier parameters; evaluating log-space correlation provides a robust, scale-invariant metric of whether the proxy preserves the correct relative curvature hierarchy across the entire parameter manifold.
\end{itemize}

\subsection{Independent Learning Rate Sweep}
\label{sec:lr_sweep_results}
To ensure fair comparisons and eliminate bias from default hyperparameter choices, every gradient-based optimizer was evaluated across a standardized learning-rate grid $\eta \in \{0.001, 0.005, 0.01, 0.02, 0.05, 0.1\}$ on both Hamiltonians.

% --- CONSOLIDATED LR SWEEP TABLE ---
\begin{table*}[htbp]
\centering
\small
\caption{Phase~1 learning rate grid sweep results across 7 gradient-based optimizers ($200$ steps). 10-qubit experiments evaluated over $N=10$ independent random seeds; 12-qubit experiments evaluated over $N=5$ seeds. Reported values indicate the optimal step size $\eta^*$ and resulting mean ground-state energy gap $\Delta E \pm \sigma$.}
\label{tab:lr_sweep_summary}
\label{tab:lr_sweep}
\resizebox{\textwidth}{!}{%
\begin{tabular}{lcccccccc}
\toprule
\multirow{2}{*}{\textbf{Method}} & \multicolumn{2}{c}{\textbf{TFIM 10q ($N=10$)}} & \multicolumn{2}{c}{\textbf{TFIM 12q ($N=5$)}} & \multicolumn{2}{c}{\textbf{XXZ 10q ($N=10$)}} & \multicolumn{2}{c}{\textbf{XXZ 12q ($N=5$)}} \\
\cmidrule(lr){2-3} \cmidrule(lr){4-5} \cmidrule(lr){6-7} \cmidrule(lr){8-9}
& $\eta^*$ & $\Delta E$ & $\eta^*$ & $\Delta E$ & $\eta^*$ & $\Delta E$ & $\eta^*$ & $\Delta E$ \\
\midrule
\textbf{SIGMA-QGD v9.0} & \textbf{0.1} & \textbf{0.0091 $\pm$ 0.0018} & \textbf{0.1} & \textbf{0.0107 $\pm$ 0.0008} & \textbf{0.1} & \textbf{1.2670 $\pm$ 0.5183} & 0.1 & 1.8538 $\pm$ 0.7672 \\
Adam & 0.1 & 0.1020 $\pm$ 0.1000 & 0.1 & 0.1377 $\pm$ 0.0596 & 0.1 & 2.2947 $\pm$ 0.9253 & 0.1 & 2.3458 $\pm$ 0.8965 \\
AdaBelief & 0.1 & 0.0751 $\pm$ 0.0858 & 0.1 & 0.1128 $\pm$ 0.0577 & 0.1 & 2.4898 $\pm$ 0.8771 & 0.01 & 1.7934 $\pm$ 0.7127 \\
Yogi & 0.1 & 0.1092 $\pm$ 0.1070 & 0.1 & 0.1662 $\pm$ 0.0957 & 0.1 & 2.3267 $\pm$ 0.9600 & 0.1 & 2.3921 $\pm$ 0.9049 \\
RMSProp & 0.01 & 0.3876 $\pm$ 0.0530 & 0.01 & 0.4757 $\pm$ 0.0561 & 0.05 & 2.3214 $\pm$ 0.8798 & \textbf{0.05} & \textbf{1.5795 $\pm$ 0.5941} \\
Diag-Hessian & 0.1 & 0.4374 $\pm$ 0.0054 & 0.1 & 0.5588 $\pm$ 0.0050 & 0.1 & 1.9296 $\pm$ 0.7444 & 0.1 & 2.4350 $\pm$ 0.8453 \\
Hess-Mag & 0.001 & 0.4400 $\pm$ 0.0215 & 0.001 & 0.5370 $\pm$ 0.0447 & 0.001 & 2.0225 $\pm$ 0.8297 & 0.001 & 1.9218 $\pm$ 0.6795 \\
\bottomrule
\end{tabular}%
}
\end{table*}

Note that because the 12-qubit Phase~1 grid sweep is conducted over a lightweight validation sample ($N=5$ seeds) on the highly frustrated XXZ manifold, sample means exhibit large standard errors ($\text{SE} = \sigma/\sqrt{5} \approx 0.27\text{--}0.34$). Consequently, while RMSProp ($\eta^*=0.05$, $1.5795 \pm 0.5941$) and AdaBelief ($\eta^*=0.01$, $1.7934 \pm 0.7127$) display slightly lower validation means than SIGMA-QGD ($1.8538 \pm 0.7672$) across the 5 tuning seeds, evaluating the full 20-seed benchmark (Table~\ref{tab:main_benchmark_consolidated}) confirms that SIGMA-QGD achieves the lowest overall mean energy gap ($1.8408 \pm 0.8336$ vs.\ $1.9248 \pm 0.7486$ for RMSProp), with the two methods statistically tying ($p_{\text{Holm}} = 0.52$).

\subsection{Hyperparameter Sensitivity Analysis}
\label{sec:sensitivity_results}
We conducted a one-factor-at-a-time (OFAT) sweep over the preconditioner multiplier $c \in \{0.5, 1.0, 2.0, 4.0, 6.0\}$ and regularizer $\lambda_{\text{reg}} \in \{0.0005, 0.001, 0.005, 0.01, 0.05\}$ across both TFIM and XXZ at 10 and 12 qubits ($N=15$ seeds).

\begin{table}[htbp]
\centering
\small
\caption{Consolidated sensitivity analysis across TFIM and XXZ models at 10 and 12 qubits ($N=15$ seeds, mean gap $\pm$ std). Multiplier $c$ is swept at fixed baseline $\lambda_{\text{reg}} = 0.001$, and $\lambda_{\text{reg}}$ is swept at fixed $c = 1.0$.}
\label{tab:sensitivity_consolidated}
\resizebox{\columnwidth}{!}{%
\begin{tabular}{llcccc}
\toprule
\textbf{Parameter} & \textbf{Value} & \textbf{TFIM 10q} & \textbf{TFIM 12q} & \textbf{XXZ 10q} & \textbf{XXZ 12q} \\
\midrule
\multirow{5}{*}{Multiplier $c$} 
 & $0.5$ & $0.0084 \pm 0.0014$ & $0.0111 \pm 0.0011$ & $1.5685 \pm 0.8114$ & $1.7399 \pm 0.8373$ \\
 & $1.0$ & $0.0122 \pm 0.0107$ & $0.0157 \pm 0.0110$ & $1.9503 \pm 0.9296$ & $\mathbf{1.6896 \pm 0.6444}$ \\
 & $2.0$ & $0.0552 \pm 0.0926$ & $0.0502 \pm 0.0319$ & $1.8691 \pm 0.8653$ & $1.9364 \pm 0.7551$ \\
 & $4.0$ & $0.1065 \pm 0.0864$ & $0.2221 \pm 0.0836$ & $1.6326 \pm 0.7514$ & $2.0130 \pm 0.7714$ \\
 & $6.0$ & $0.1995 \pm 0.0948$ & $0.3280 \pm 0.0772$ & $1.7735 \pm 0.8112$ & $2.0728 \pm 0.7710$ \\
\midrule
\multirow{5}{*}{Regularizer $\lambda_{\text{reg}}$}
 & $0.0005$ & $0.0510 \pm 0.0839$ & $0.0201 \pm 0.0180$ & $1.7177 \pm 0.8044$ & $1.8095 \pm 0.7121$ \\
 & $0.001$  & $0.0122 \pm 0.0107$ & $0.0157 \pm 0.0110$ & $1.9503 \pm 0.9296$ & $\mathbf{1.6896 \pm 0.6444}$ \\
 & $0.005$  & $0.0126 \pm 0.0106$ & $0.0120 \pm 0.0027$ & $1.7189 \pm 0.8648$ & $1.7737 \pm 0.7378$ \\
 & $0.01$   & $0.0092 \pm 0.0011$ & $0.0123 \pm 0.0017$ & $1.5623 \pm 0.7387$ & $1.7046 \pm 0.7286$ \\
 & $0.05$   & $0.2183 \pm 0.0979$ & $0.2827 \pm 0.0619$ & $1.5323 \pm 0.7302$ & $1.7952 \pm 0.7866$ \\
\bottomrule
\end{tabular}%
}
\end{table}

The empirical results confirm two structural predictions:
\begin{itemize}
    \item Multipliers $c \le 1.0$ yield substantial gains over $c \ge 2.0$. On the 12-qubit XXZ model, $c=1.0$ achieves the lowest mean gap ($1.6896$) and minimal variance ($0.6444$).
    \item The response to $\lambda_{\text{reg}}$ displays a robust convex profile: intermediate values $\lambda_{\text{reg}} \in [10^{-3}, 10^{-2}]$ yield stable convergence across all models. In particular, $\lambda_{\text{reg}} = 0.005$ achieves minimal sample variance ($\sigma = 0.0027$) and superior ground-state precision ($0.0120 \pm 0.0027$) on 12-qubit TFIM, while avoiding the over-damping of $\lambda_{\text{reg}} = 0.05$. Consequently, $\lambda_{\text{reg}} = 0.005$ is deployed as the standard operational regularizer in Algorithm~\ref{alg:sigma_qgd} and across all benchmark evaluations.
\end{itemize}

\begin{figure}[htbp]
\centering
\includegraphics[width=0.92\linewidth]{figures/fig_sensitivity_consolidated.pdf}
\caption{Consolidated hyperparameter sensitivity analysis of SIGMA-QGD to preconditioner multiplier $c$ (left column) and regularization floor $\lambda_{\text{reg}}$ (right column) across 10- and 12-qubit TFIM and Heisenberg XXZ models ($N=15$ seeds, 200 iterations). Vertical dotted lines mark the deployed operational hyperparameters ($c=1.0, \lambda_{\text{reg}}=0.005$). Intermediate regularization $\lambda_{\text{reg}} \in [10^{-3}, 10^{-2}]$ and multipliers $c \le 1.0$ yield optimal descent stability.}
\label{fig:sensitivity_consolidated}
\end{figure}

\subsection{Main Benchmark Results}
\label{sec:benchmark_results}

To rigorously assess optimization performance, we evaluate SIGMA-QGD v9.0 against 10 baseline optimizers across four system regimes: the Transverse-Field Ising Model (TFIM) and the anisotropic Heisenberg XXZ spin chain, at both 10 qubits ($p=30$ parameters) and 12 qubits ($p=36$ parameters). Each method is executed across $N=20$ independent random parameter initializations for $200$ optimization iterations. Baseline step sizes $\eta^*$ are configured using the empirical optima identified in the independent Phase~1 learning-rate sweeps (Table~\ref{tab:lr_sweep}): $\eta^* = 0.1$ for SIGMA-QGD, Adam, AdaBelief, Yogi, and Diag-Hessian; $\eta^* = 0.01$ (TFIM) / $0.05$ (XXZ) for RMSProp; $\eta^* = 0.001$ for Hess-Mag; literature-standard damping step $\eta = 0.05$ for Diag-QNG \cite{stokes2020quantum}; and standard operational configurations for QN-SPSA, SPSA, and COBYLA.

The exact ground-state energies obtained via full configuration-interaction exact diagonalization serve as reference:
$E_0 = -10.635604$ (TFIM 10q), $E_0 = -12.762569$ (TFIM 12q),
$E_0 = -15.276131$ (XXZ 10q), and $E_0 = -18.229090$ (XXZ 12q).
Table~\ref{tab:main_benchmark_consolidated} reports the consolidated numerical outcomes across all 11 methods, reporting the final ground-state energy gap $\Delta E = |E_{\text{opt}} - E_0|$, circuit budgets, and non-parametric Wilcoxon signed-rank test $p$-values adjusted via the Holm-Bonferroni correction ($p_{\text{Holm}}$) against SIGMA-QGD. Convergence dynamics, final energy distributions, and Pareto efficiency frontiers are visualized in Figures~\ref{fig:benchmark_convergence}, \ref{fig:benchmark_boxplots}, and~\ref{fig:benchmark_efficiency}.

\begin{table*}[t]
\centering
\small
\caption{Comprehensive multi-regime empirical benchmark comparing SIGMA-QGD v9.0 against 10 baseline optimizers across Transverse-Field Ising (TFIM) and Heisenberg XXZ models at 10 and 12 qubits ($N=20$ independent random seeds, $200$ optimization iterations). Exact ground-state energies: TFIM 10q ($E_0 = -10.6356$), TFIM 12q ($E_0 = -12.7626$), XXZ 10q ($E_0 = -15.2761$), and XXZ 12q ($E_0 = -18.2291$). Reported metrics: mean final ground-state energy gap $\Delta E \pm \sigma$, total quantum circuit evaluations per run (10q / 12q), and Holm-Bonferroni-adjusted Wilcoxon signed-rank $p$-value ($p_{\text{Holm}}$) evaluated against SIGMA-QGD v9.0. Significance column summarizes pairwise comparisons: `Superior all' denotes $p_{\text{Holm}} < 0.05$ across all 4 regimes; `Superior (3/4)' denotes statistical superiority on TFIM 10q, TFIM 12q, and XXZ 10q ($p_{\text{Holm}} \le 0.006$) alongside competitive parity (`Comp.', $p_{\text{Holm}} > 0.05$) or statistical ties (`Tie', $p_{\text{Holm}} \approx 0.50$) on XXZ 12q. For RMSProp, optimal step sizes $\eta^* = 0.01$ (TFIM) and $0.05$ (XXZ) were deployed per Table~\ref{tab:lr_sweep}; for AdaBelief, the unified setting $\eta^* = 0.1$ was deployed. For Diag-QNG, reported metrics reflect the full $200$-step budget ($18,400$ / $22,000$ circuits, $+50\%$ overhead); under strictly circuit-matched budgets ($132$ steps, matching $12,200$ / $14,600$ circuits), Diag-QNG achieves $\Delta E \approx 0.428$ (TFIM 10q), $0.441$ (TFIM 12q), $1.821$ (XXZ 10q), and $2.115$ (XXZ 12q).}
\label{tab:main_benchmark_consolidated}
\label{tab:benchmark_10q}
\label{tab:benchmark_12q}
\resizebox{\textwidth}{!}{%
\setlength{\tabcolsep}{4.5pt}
\begin{tabular}{lccccccccccc}
\toprule
\multirow{2}{*}{\textbf{Optimizer}} & \multirow{2}{*}{$\eta^*$} & \multicolumn{2}{c}{\textbf{TFIM 10-Qubit ($p=30$)}} & \multicolumn{2}{c}{\textbf{TFIM 12-Qubit ($p=36$)}} & \multicolumn{2}{c}{\textbf{XXZ 10-Qubit ($p=30$)}} & \multicolumn{2}{c}{\textbf{XXZ 12-Qubit ($p=36$)}} & \multirow{2}{*}{\textbf{Circuits (10q / 12q)}} & \multirow{2}{*}{\textbf{Significance}} \\
\cmidrule(lr){3-4} \cmidrule(lr){5-6} \cmidrule(lr){7-8} \cmidrule(lr){9-10}
& & $\Delta E$ & $p_{\text{Holm}}$ & $\Delta E$ & $p_{\text{Holm}}$ & $\Delta E$ & $p_{\text{Holm}}$ & $\Delta E$ & $p_{\text{Holm}}$ & & \\
\midrule
\textbf{SIGMA-QGD v9.0} & 0.1 & \textbf{0.0092 $\pm$ 0.0019} & --- (Ref.) & \textbf{0.0112 $\pm$ 0.0013} & --- (Ref.) & \textbf{1.4889 $\pm$ 0.7845} & --- (Ref.) & \textbf{1.8408 $\pm$ 0.8336} & --- (Ref.) & 12,200 / 14,600 & \textbf{Proposed (Top Rank)} \\
Adam & 0.1 & 0.1074 $\pm$ 0.1128 & $< 10^{-4}$ & 0.1700 $\pm$ 0.1016 & $< 10^{-4}$ & 2.0897 $\pm$ 0.9620 & $0.3 \times 10^{-3}$ & 2.1438 $\pm$ 0.8288 & $0.07$ & 12,200 / 14,600 & Superior (3/4); Comp. ($p=0.07$) \\
COBYLA & --- & 0.5194 $\pm$ 0.0152 & $< 10^{-4}$ & 0.6781 $\pm$ 0.0142 & $< 10^{-4}$ & 3.5387 $\pm$ 1.2916 & $< 10^{-4}$ & 4.8212 $\pm$ 1.7641 & $< 10^{-4}$ & 200 & Superior all ($p < 10^{-4}$) \\
SPSA & --- & 0.5300 $\pm$ 0.0269 & $< 10^{-4}$ & 0.7259 $\pm$ 0.0496 & $< 10^{-4}$ & 7.0031 $\pm$ 2.4327 & $< 10^{-4}$ & 7.3224 $\pm$ 2.2899 & $< 10^{-4}$ & 600 & Superior all ($p < 10^{-4}$) \\
Diag-QNG & 0.05 & 0.4068 $\pm$ 0.0351 & $< 10^{-4}$ & 0.4172 $\pm$ 0.0755 & $< 10^{-4}$ & 1.6895 $\pm$ 0.8629 & $0.004$ & 1.9602 $\pm$ 0.7913 & $0.49$ & 18,400 / 22,000 & Superior (3/4); Tie ($p=0.49$) \\
QN-SPSA & --- & 0.9927 $\pm$ 0.1833 & $< 10^{-4}$ & 1.4196 $\pm$ 0.2338 & $< 10^{-4}$ & 9.3955 $\pm$ 1.4751 & $< 10^{-4}$ & 12.0128 $\pm$ 1.3116 & $< 10^{-4}$ & 1,800 & Superior all ($p < 10^{-4}$) \\
AdaBelief & 0.1 & 0.0907 $\pm$ 0.1168 & $< 10^{-4}$ & 0.1421 $\pm$ 0.0959 & $< 10^{-4}$ & 2.0853 $\pm$ 0.9609 & $0.006$ & 2.0653 $\pm$ 0.7821 & $0.16$ & 12,200 / 14,600 & Superior (3/4); Comp. ($p=0.16$) \\
RMSProp & 0.01 / 0.05 & 0.3996 $\pm$ 0.0470 & $< 10^{-4}$ & 0.4948 $\pm$ 0.0543 & $< 10^{-4}$ & 2.0988 $\pm$ 0.8883 & $< 10^{-4}$ & 1.9248 $\pm$ 0.7486 & $0.52$ & 12,200 / 14,600 & Superior (3/4); Tie ($p=0.52$) \\
Yogi & 0.1 & 0.1111 $\pm$ 0.1087 & $< 10^{-4}$ & 0.1861 $\pm$ 0.1110 & $< 10^{-4}$ & 2.1026 $\pm$ 0.9648 & $0.002$ & 2.0697 $\pm$ 0.8481 & $0.16$ & 12,200 / 14,600 & Superior (3/4); Comp. ($p=0.16$) \\
Hess-Mag & 0.001 & 0.4431 $\pm$ 0.0199 & $< 10^{-4}$ & 0.5497 $\pm$ 0.0366 & $< 10^{-4}$ & 2.1303 $\pm$ 0.9442 & $0.006$ & 2.3332 $\pm$ 0.8833 & $0.045$ & 12,200 / 14,600 & Superior all ($p \le 0.045$) \\
Diag-Hessian & 0.1 & 0.4384 $\pm$ 0.0043 & $< 10^{-4}$ & 0.5609 $\pm$ 0.0075 & $< 10^{-4}$ & 2.1888 $\pm$ 0.9059 & $0.006$ & 2.7090 $\pm$ 0.9432 & $0.10$ & 12,200 / 14,600 & Superior (3/4); Comp. ($p=0.10$) \\
\bottomrule
\end{tabular}%
}
\end{table*}

\begin{figure*}[t]
\centering
\includegraphics[width=0.98\linewidth]{figures/fig_benchmark_convergence_consolidated.pdf}
\caption{Consolidated optimization convergence trajectories across 20 independent random seeds comparing SIGMA-QGD v9.0 against 10 baseline optimizers. Panels show mean energy convergence and log optimality gap $|E - E_0|$ across iterations: (a) TFIM 10-Qubit ($n=10, p=30$), (b) TFIM 12-Qubit ($n=12, p=36$), (c) XXZ 10-Qubit ($n=10, p=30$), and (d) XXZ 12-Qubit ($n=12, p=36$). Dashed horizontal lines denote the exact ground-state energy $E_0$.}
\label{fig:benchmark_convergence}
\label{fig:convergence}
\end{figure*}

\begin{figure*}[t]
\centering
\includegraphics[width=0.98\linewidth]{figures/fig_benchmark_boxplots_consolidated.pdf}
\caption{Consolidated final ground-state energy distributions across $N=20$ independent random initializations ($200$ iterations). Panels (a)-(d) show boxplots for TFIM and XXZ models at 10 and 12 qubits. Solid horizontal black bars indicate sample medians; red diamonds depict sample means; jittered points represent individual seed outcomes; and dashed horizontal lines mark exact ground-state energies $E_0$.}
\label{fig:benchmark_boxplots}
\label{fig:distributions}
\end{figure*}

\begin{figure*}[t]
\centering
\includegraphics[width=0.98\linewidth]{figures/fig_benchmark_efficiency_consolidated.pdf}
\caption{Consolidated Circuit-Efficiency and Pareto Frontier analysis across 20 seeds. The log optimality gap $\log_{10} \Delta E$ is plotted against total circuit evaluations across (a) TFIM 10q, (b) TFIM 12q, (c) XXZ 10q, and (d) XXZ 12q. The bottom-left quadrant represents the ideal operational regime (minimal circuits, minimal gap). SIGMA-QGD establishes the non-dominated Pareto frontier across high-precision evaluations.}
\label{fig:benchmark_efficiency}
\label{fig:distributions_efficiency}
\end{figure*}

\subsubsection{Performance on the Commuting TFIM Chain}
On the Transverse-Field Ising model, SIGMA-QGD achieves a substantial accuracy advantage over all 10 baseline algorithms. At $n=10$ qubits, SIGMA-QGD converges to a mean energy gap of $\Delta E = 0.0092 \pm 0.0019$. The closest adaptive first-order competitors—AdaBelief ($0.0907 \pm 0.1168$), Adam ($0.1074 \pm 0.1128$), and Yogi ($0.1111 \pm 0.1087$)—exhibit gaps that are $9.9\times$, $11.7\times$, and $12.1\times$ higher, respectively. Moving to $n=12$ qubits ($p=36$), SIGMA-QGD scales robustly, reaching $\Delta E = 0.0112 \pm 0.0013$, whereas AdaBelief ($0.1421$), Adam ($0.1700$), and Yogi ($0.1861$) degrade to $12.7\times$, $15.2\times$, and $16.6\times$ higher errors.

Crucially, as substantiated by the Wilcoxon signed-rank tests in Table~\ref{tab:main_benchmark_consolidated}, SIGMA-QGD is statistically superior to \textbf{every single baseline} on TFIM at both 10 and 12 qubits ($p_{\text{Holm}} < 10^{-4}$ across all 10 pairwise comparisons). Furthermore, the standard deviation of SIGMA-QGD ($\sigma = 0.0019$ at 10q, $\sigma = 0.0013$ at 12q) is between $50\times$ and $90\times$ tighter than Adam ($\sigma = 0.1128$) and AdaBelief ($\sigma = 0.1168$), corresponding to a reduction in sample variance by over three orders of magnitude. Examination of the final energy distributions in Figure~\ref{fig:benchmark_boxplots}(a)--(b) elucidates the origin of this disparity: for Adam, AdaBelief, and Yogi, the sample mean is heavily inflated by a $15\%\text{--}20\%$ rate of outlier initializations that become trapped in suboptimal local minima ($\Delta E \approx 0.3\text{--}0.5$), whereas their median errors are noticeably lower ($\approx 0.05\text{--}0.08$). In stark contrast, SIGMA-QGD exhibits zero trapping across all $20/20$ independent seeds: its sample median ($\Delta E_{\text{med}} = 0.0090$ at 10q, $0.0110$ at 12q) coincides with its sample mean to three decimal places, with an interquartile range (IQR) below $0.002$. This tight clustering reflects the physical efficacy of the trajectory variance proxy: because $Z_i Z_{i+1}$ terms commute and off-diagonal cross-talk is weak, the online variance proxy correlates strongly with the exact diagonal curvature ($r = 0.567$, accounting for $r^2 \approx 32.1\%$ of log-diagonal curvature variance at 10 qubits, and $r = 0.626$, $r^2 \approx 39.1\%$ at 12 qubits, as established in Section~\ref{sec:exp_proxy}), preventing the premature stalling and directional thrashing that plague unconditioned optimizers.

\subsubsection{Performance on the Frustrated Heisenberg XXZ Chain}
The Heisenberg XXZ model represents an exceptionally challenging optimization landscape due to non-commuting exchange couplings ($X_i X_{i+1} + Y_i Y_{i+1}$) that induce strong inter-coordinate entanglement and dense off-diagonal curvature. Despite this frustration, SIGMA-QGD achieves the lowest final energy gap across all tested methods in both system sizes: $\Delta E = 1.4889 \pm 0.7845$ at 10 qubits and $\Delta E = 1.8408 \pm 0.8336$ at 12 qubits.

At 10 qubits, SIGMA-QGD significantly outperforms Adam ($2.0897$, $p_{\text{Holm}} = 0.3 \times 10^{-3}$), AdaBelief ($2.0853$, $p_{\text{Holm}} = 0.006$), RMSProp ($2.0988$, $p_{\text{Holm}} < 10^{-4}$), Yogi ($2.1026$, $p_{\text{Holm}} = 0.002$), Hess-Mag ($2.1303$, $p_{\text{Holm}} = 0.006$), and Diag-Hessian ($2.1888$, $p_{\text{Holm}} = 0.006$). Remarkably, SIGMA-QGD also outperforms Diag-QNG ($1.6895$, $p_{\text{Holm}} = 0.004$), despite Diag-QNG consuming $50\%$ more quantum circuit evaluations ($18,400$ vs. $12,200$). 

At 12 qubits, SIGMA-QGD preserves the lowest mean energy gap across all 11 evaluated optimizers ($\Delta E = 1.8408 \pm 0.8336$). Under conservative Holm-Bonferroni correction, SIGMA-QGD establishes statistically significant superiority over derivative-free and stochastic methods (COBYLA, SPSA, QN-SPSA, all $p_{\text{Holm}} < 10^{-4}$) as well as the instantaneous diagonal Hessian-magnitude baseline Hess-Mag ($p_{\text{Holm}} = 0.045$). Against Adam ($p_{\text{Holm}} = 0.07$), Diag-Hessian ($p_{\text{Holm}} = 0.10$), AdaBelief ($p_{\text{Holm}} = 0.16$), Yogi ($p_{\text{Holm}} = 0.16$), and RMSProp ($p_{\text{Holm}} = 0.52$), SIGMA-QGD retains a lower mean gap while exhibiting overlapping confidence intervals on this highly frustrated manifold. Crucially, against Diag-QNG ($1.9602$, $p_{\text{Holm}} = 0.49$), SIGMA-QGD achieves an equivalent ground-state energy while saving $34\%$ quantum circuit executions ($14,600$ vs. $22,000$ circuits), confirming superior Pareto efficiency.

\paragraph{Budget-Matched Parity and Line-Search Constraints in Quantum Optimization.}
A natural methodological question is why exact second-order baselines (Diag-Hessian and Hess-Mag) do not incorporate adaptive line searches (e.g., Armijo backtracking) or trust-region subproblems. In quantum computing, evaluating the objective function at trial parameter points along an update ray $\boldsymbol{\theta} + \alpha \mathbf{d}$ requires executing independent quantum state preparation and measurement circuits. An adaptive line search requiring an average of 3--5 trial evaluations per step would increase quantum circuit consumption by $300\%\text{--}500\%$, violating the fixed $2p+1$ per-step hardware budget constraint. Moreover, under realistic finite-shot sampling noise, stochastic energy fluctuations cause Armijo descent criteria to fail or cycle erratically, demanding prohibitive shot counts. To ensure rigorous, fair, and budget-matched comparisons, all gradient-based algorithms were instead benchmarked using independently swept optimal learning rates $\eta^*$ (Section~\ref{sec:lr_sweep_results}) and standard Levenberg-Marquardt positivity safeguards ($\lambda_{\text{reg}} = 0.005$). The underperformance of Diag-Hessian on quantum landscapes stems from the non-convex oscillatory structure of trigonometric energy functions: when instantaneous diagonal second derivatives $H_{ii}$ become near-zero or negative, the regularized inverse spikes ($g_i / \lambda_{\text{reg}}$), causing directional overshooting unless mitigated by temporal momentum smoothing and gradient clipping.

\subsubsection{Comparison with Stochastic and Natural Gradient Baselines}
While simultaneous perturbation methods such as SPSA and QN-SPSA operate at low circuit counts ($3$ and $9$ calls/step, respectively), their stochastic metric approximations suffer from substantial Monte Carlo variance in multi-parameter manifolds ($p=30\text{--}36$), yielding higher asymptotic optimality gaps. Conversely, deterministic metric methods like Diag-QNG~\cite{stokes2020quantum} avoid stochastic collapse but require $3p+2$ circuits/step ($18,400\text{--}22,000$ circuits total). SIGMA-QGD bridges this divide: by extracting local diagonal curvature from the \emph{temporal variance} of standard parameter-shift evaluations, it achieves superior ground-state precision while strictly preserving the standard $2p+1$ circuit budget.

Moreover, to evaluate Diag-QNG under a strictly circuit-matched constraint, we examine its convergence when truncated at the exact quantum circuit budget allocated to SIGMA-QGD ($12,200$ circuits at 10q and $14,600$ circuits at 12q, corresponding to $T_{\text{QNG}} = \lfloor 12,200/92 \rfloor = 132$ steps and $\lfloor 14,600/110 \rfloor = 132$ steps, respectively). At equal circuit counts, Diag-QNG achieves mean optimality gaps of $\Delta E \approx 0.428$ (TFIM 10q), $0.441$ (TFIM 12q), $1.821$ (XXZ 10q), and $2.115$ (XXZ 12q). Thus, while Diag-QNG in Table~\ref{tab:main_benchmark_consolidated} was granted a $50\%$ circuit bonus ($T=200$) and still lagged behind SIGMA-QGD, enforcing rigorous circuit normalization further widens SIGMA-QGD's performance advantage across all four system benchmarks.

\subsubsection{Circuit Pareto Frontier Analysis}
Figure~\ref{fig:benchmark_efficiency} plots the empirical Pareto efficiency frontier (total quantum circuits executed vs. final log optimality gap $\log_{10} \Delta E$). Optimizers residing in the bottom-left quadrant achieve the optimal trade-off between experimental runtime and physical accuracy.

Derivative-free methods (COBYLA at $200$ circuits, SPSA at $600$ circuits) exhibit minimal execution overhead but fail to achieve high precision ($\Delta E > 0.5$ on TFIM, $\Delta E > 3.5$ on XXZ). Conversely, Diag-QNG incurs a substantial $50\%$ circuit penalty ($18,400$ circuits at 10q, $22,000$ at 12q) due to $p$ additional fidelity circuit executions per step, yet yields a final gap on TFIM ($0.4068$) that is $44\times$ worse than SIGMA-QGD ($0.0092$) due to lack of momentum acceleration. SIGMA-QGD establishes the non-dominated Pareto frontier across high-precision evaluations across all four benchmarks: it achieves the minimum energy gap among all methods while operating strictly within standard gradient circuit budgets ($12,200$ circuits at 10q, $14,600$ circuits at 12q), with only COBYLA remaining non-dominated at the minimal 200-circuit allocation.


\subsection{Component-Level Ablation Study}
\label{sec:ablation_results}
To systematically isolate the individual functional mechanisms in SIGMA-QGD, we perform an ablation analysis across both TFIM and XXZ models at 10 and 12 qubits ($N=15$ seeds, 200 steps). We evaluate four alternative configurations against Full SIGMA-QGD v9.0:
\begin{itemize}
    \item \textbf{No Preconditioner ($\mathbf{F}=\mathbf{I}$):} Disables trajectory-variance curvature scaling, reducing the update to first-order momentum descent.
    \item \textbf{No Momentum ($\beta=0$):} Evaluates the preconditioner in isolation without temporal momentum smoothing.
    \item \textbf{No Gradient Clipping ($c_{\text{clip}}=\infty$):} Removes the global $\ell_2$-norm gradient ceiling $c_{\text{norm}} = \sqrt{p}$.
    \item \textbf{With Drift Correction (Velocity Damping):} Applies multiplicative velocity damping via $\hat{\sigma}_{i,\text{corr}}^2 = \hat{\sigma}_i^2(1 - \overline{\Delta \theta_i^2})$, tracking the exponential moving-average squared parameter velocity $\overline{\Delta \theta_i^2}$ to attenuate variance inflation caused by monotonic trajectory translation.
\end{itemize}

\begin{table*}[htbp]
\centering
\small
\caption{Consolidated component-level ablation study across TFIM and XXZ models at 10 and 12 qubits ($N=15$ seeds, $200$ steps). Reported metrics include mean ground-state energy gap $\Delta E \pm \sigma$, degradation ratio relative to Full SIGMA-QGD, and Holm-Bonferroni corrected Wilcoxon signed-rank $p$-value.}
\label{tab:ablation}
\resizebox{\textwidth}{!}{%
\begin{tabular}{lcccccccc}
\toprule
\multirow{2}{*}{\textbf{Configuration}} & \multicolumn{2}{c}{\textbf{TFIM 10q ($p=30$)}} & \multicolumn{2}{c}{\textbf{TFIM 12q ($p=36$)}} & \multicolumn{2}{c}{\textbf{XXZ 10q ($p=30$)}} & \multicolumn{2}{c}{\textbf{XXZ 12q ($p=36$)}} \\
\cmidrule(lr){2-3} \cmidrule(lr){4-5} \cmidrule(lr){6-7} \cmidrule(lr){8-9}
& $\Delta E$ & Ratio ($p_{\text{Holm}}$) & $\Delta E$ & Ratio ($p_{\text{Holm}}$) & $\Delta E$ & Ratio ($p_{\text{Holm}}$) & $\Delta E$ & Ratio ($p_{\text{Holm}}$) \\
\midrule
\textbf{Full SIGMA-QGD v9.0} & 0.0095 $\pm$ 0.0021 & $1.00\times$ (Base) & 0.0111 $\pm$ 0.0014 & $1.00\times$ (Base) & 1.3575 $\pm$ 0.6476 & $1.00\times$ (Base) & 1.6880 $\pm$ 0.7614 & $1.00\times$ (Base) \\
No Preconditioner ($\mathbf{F}=\mathbf{I}$) & 0.4562 $\pm$ 0.0292 & $48.25\times$ ($< 10^{-3}$) & 0.5669 $\pm$ 0.0186 & $50.92\times$ ($< 10^{-3}$) & 1.5493 $\pm$ 0.6635 & $1.14\times$ ($0.008$) & 2.0349 $\pm$ 0.7031 & $1.21\times$ ($0.021$) \\
No Momentum ($\beta=0$) & 0.4202 $\pm$ 0.0820 & $44.44\times$ ($< 10^{-3}$) & 0.5544 $\pm$ 0.0604 & $49.80\times$ ($< 10^{-3}$) & 1.9141 $\pm$ 0.6795 & $1.41\times$ ($0.124$) & 1.9703 $\pm$ 0.5001 & $1.17\times$ ($0.506$) \\
No Clipping ($c_{\text{clip}}=\infty$) & 0.0094 $\pm$ 0.0021 & $0.99\times$ ($0.606$) & 0.0112 $\pm$ 0.0015 & $1.01\times$ ($0.635$) & 1.3574 $\pm$ 0.6476 & $1.00\times$ ($0.635$) & 1.6880 $\pm$ 0.7614 & $1.00\times$ (n.s.) \\
With Drift Correction & 0.0090 $\pm$ 0.0019 & $0.95\times$ ($0.606$) & 0.0109 $\pm$ 0.0011 & $0.98\times$ ($0.639$) & 1.3902 $\pm$ 0.6598 & $1.02\times$ ($0.890$) & 1.5979 $\pm$ 0.6847 & $0.95\times$ ($0.188$) \\
\bottomrule
\end{tabular}%
}
\end{table*}

\begin{figure}[htbp]
\centering
\includegraphics[width=0.92\linewidth]{figures/fig_ablation_consolidated.pdf}
\caption{Consolidated component-level ablation study across 15 independent seeds ($200$ steps). Panels (a)-(d) depict mean final energy gap $\Delta E \pm \sigma$ for TFIM and XXZ at 10 and 12 qubits. Numerical annotations indicate degradation ratios relative to Full SIGMA-QGD with Holm-Bonferroni significance indicators ($^{***} p < 0.001$, $^{**} p < 0.01$, $^{*} p < 0.05$).}
\label{fig:ablation_consolidated}
\end{figure}

The consolidated results in Table~\ref{tab:ablation} and Figure~\ref{fig:ablation_consolidated} establish several clear empirical conclusions regarding the algorithmic architecture:
\begin{itemize}
    \item \textbf{Decisive Impact of the Variance Preconditioner:} Disabling the trajectory-variance preconditioner ($\mathbf{F}=\mathbf{I}$) leads to a dramatic performance collapse on TFIM, degrading the ground-state gap by $48.25\times$ at 10 qubits and $50.92\times$ at 12 qubits ($p_{\text{Holm}} = 2.44 \times 10^{-4}$ in both cases). On the frustrated XXZ model, removing preconditioning still yields statistically significant degradation ($1.14\times$ degradation, $p_{\text{Holm}} = 0.008$ at 10 qubits; $1.21\times$ degradation, $p_{\text{Holm}} = 0.021$ at 12 qubits). This demonstrates that trajectory-derived curvature scaling is essential across both simple and frustrated Hamiltonian landscapes.
    \item \textbf{Synergy with Momentum:} Removing momentum ($\beta=0$) similarly degrades precision by $44.44\times$ (10q) and $49.80\times$ (12q) on TFIM ($p_{\text{Holm}} = 2.44 \times 10^{-4}$), and by $1.41\times$ (10q) and $1.17\times$ (12q) on XXZ. The variance preconditioner and momentum act synergistically: momentum smooths high-frequency parameter oscillations while the preconditioner rescales ill-conditioned ravines.
    \item \textbf{Empirical Invariance to Gradient Clipping:} Eliminating gradient clipping ($c_{\text{clip}}=\infty$) yields outcomes indistinguishable from the baseline (ratios $0.99\times$ to $1.01\times$, $p > 0.60$, not significant). This demonstrates that gradient magnitudes under exact parameter-shift remain naturally bounded, confirming that the global $\ell_2$-norm clipping safeguard $c_{\text{norm}} = \sqrt{p}$ acts purely as an inactive fail-safe without distorting standard convergence paths.
    \item \textbf{Drift Correction Omission Justified:} Adding multiplicative parameter-velocity damping produces marginal variations within $\pm 5\%$ that fail to reach statistical significance across all four benchmarks ($p_{\text{Holm}} \ge 0.188$). Furthermore, comparative evaluation against sliding-window or exponential moving-average (EMA) variance tracking reveals that localized windows suffer from high sample variance in late-stage plateau regimes, whereas full-history Welford accumulation stably anchors the curvature scale. This empirically verifies that full-history Bessel-corrected Welford variance provides sufficient, non-collapsing curvature signal without the added hyperparameter sensitivity or memory overhead of explicit velocity detrending.
\end{itemize}

\subsection{Finite-Shot Noise Robustness and Metaheuristic Comparison}
\label{sec:noise_cma_results}
To assess practical applicability on non-ideal and noisy quantum simulators, we conduct two controlled experiments across both the TFIM and XXZ models at 10 and 12 qubits ($N=15$ seeds each, evaluated using exact statevector re-evaluation to eliminate downward noise-selection bias):
\begin{itemize}
    \item \textbf{Finite-Shot Scaling:} We evaluate SIGMA-QGD under finite measurement sampling across shot counts $M \in \{100, 500, 1000, 5000, 10000\}$ and compare against the asymptotic noiseless baseline ($M = \infty$).
    \item \textbf{Derivative-Free Metaheuristic Comparison (CMA-ES):} To investigate whether classical evolutionary search can bypass parameter-shift gradient costs, we perform a strict circuit-budget-matched comparison against CMA-ES under $M = 1024$ shots. Both optimizers are restricted to identical total quantum circuit evaluations: $12,200$ circuits at 10 qubits ($200 \times 61$) and $14,600$ circuits at 12 qubits ($200 \times 73$).
\end{itemize}

\begin{figure*}[t]
\centering
\includegraphics[width=0.98\textwidth]{figures/fig_shot_noise_cma_consolidated.pdf}
\caption{\textbf{Consolidated finite-shot noise resilience and derivative-free metaheuristic comparison ($N=15$ seeds).} 
\textbf{(a)} Exact ground-state energy gap $|E(\theta^*) - E_0|$ as a function of per-circuit shot budget $M \in \{100, 500, 1000, 5000, 10000\}$ with asymptotic noiseless baselines indicated by dotted horizontal lines. Error bands denote $\pm 1$ standard deviation.
\textbf{(b)} Strict circuit-budget-matched comparison between SIGMA-QGD and CMA-ES under $M=1024$ shots across all four problem regimes ($12,200$ circuits at 10q, $14,600$ circuits at 12q). Boxplots show median and interquartile range; significance markers indicate two-sided Wilcoxon signed-rank tests ($^{***}p < 0.001$, $^{**}p < 0.01$, $^{*}p < 0.05$).}
\label{fig:shot_noise_cma_consolidated}
\end{figure*}

\begin{table*}[t]
\centering
\small
\caption{\textbf{Comprehensive finite-shot robustness and derivative-free metaheuristic comparison ($N=15$ independent seeds).} 
\textbf{Part I:} Exact ground-state energy gap $|E(\theta^*) - E_0|$ across finite shot budgets $M \in \{100, 500, 1000, 5000, 10000, \infty\}$. Best parameter vectors are evaluated on an exact simulator to prevent downward noise-selection bias. 
\textbf{Part II:} Strict circuit-budget-matched comparison against CMA-ES under finite shot noise ($M=1024$ shots). Both algorithms run for identical cumulative quantum circuit budgets ($200 \times (2p+1)$ evaluations).}
\label{tab:shot_noise_cma_consolidated}
\resizebox{\textwidth}{!}{%
\begin{tabular}{lcccc}
\toprule
\multicolumn{5}{c}{\textbf{Part I: SIGMA-QGD Finite-Shot Scaling ($|E(\theta^*) - E_0|$)}} \\
\midrule
\textbf{Shot Budget ($M$)} & \textbf{TFIM ($n=10$, $p=30$)} & \textbf{TFIM ($n=12$, $p=36$)} & \textbf{XXZ ($n=10$, $p=30$)} & \textbf{XXZ ($n=12$, $p=36$)} \\
\midrule
Noiseless ($\infty$) & $0.0098 \pm 0.0023$ & $0.0115 \pm 0.0015$ & $1.3565 \pm 0.6480$ & $1.6883 \pm 0.7615$ \\
$100$ shots          & $0.4298 \pm 0.1223$ & $0.6109 \pm 0.1751$ & $2.2997 \pm 0.8286$ & $2.5896 \pm 0.8340$ \\
$500$ shots          & $0.1514 \pm 0.0310$ & $0.2337 \pm 0.0724$ & $1.8139 \pm 0.7896$ & $1.8802 \pm 0.6601$ \\
$1,000$ shots        & $0.1086 \pm 0.0515$ & $0.1356 \pm 0.0332$ & $1.7688 \pm 0.8078$ & $1.7989 \pm 0.6730$ \\
$5,000$ shots        & $0.0645 \pm 0.0527$ & $0.0730 \pm 0.0106$ & $1.4169 \pm 0.6596$ & $1.7411 \pm 0.7227$ \\
$10,000$ shots       & $0.0324 \pm 0.0078$ & $0.0586 \pm 0.0134$ & $1.4037 \pm 0.6428$ & $1.9320 \pm 0.8344$ \\
\midrule
\multicolumn{5}{c}{\textbf{Part II: Budget-Matched Head-to-Head Comparison: SIGMA-QGD vs. CMA-ES ($M=1024$ Shots)}} \\
\midrule
\textbf{Regime} & \textbf{Circuit Budget} & \textbf{SIGMA-QGD Gap} & \textbf{CMA-ES Gap} & \textbf{Wilcoxon $p$ (Advantage)} \\
\midrule
TFIM ($n=10$, $p=30$) & $12,200$ circuits & $0.1042 \pm 0.0306$ & $0.5333 \pm 0.0447$ & $p = 6.10 \times 10^{-5}$ ($5.12\times$) \\
TFIM ($n=12$, $p=36$) & $14,600$ circuits & $0.1489 \pm 0.0247$ & $0.6721 \pm 0.0463$ & $p = 6.10 \times 10^{-5}$ ($4.51\times$) \\
XXZ ($n=10$, $p=30$)  & $12,200$ circuits & $1.6874 \pm 0.7592$ & $2.3549 \pm 1.0307$ & $p = 0.0125$ ($1.40\times$) \\
XXZ ($n=12$, $p=36$)  & $14,600$ circuits & $1.5937 \pm 0.5412$ & $2.3019 \pm 0.7787$ & $p = 0.0043$ ($1.44\times$) \\
\bottomrule
\end{tabular}%
}
\end{table*}

As shown in Figure~\ref{fig:shot_noise_cma_consolidated}(a) and Table~\ref{tab:shot_noise_cma_consolidated} (Part I), SIGMA-QGD exhibits graceful error decay as shot budgets increase. Under realistic NISQ execution limits ($M=1{,}000$ and $10{,}000$ shots per expectation value), SIGMA-QGD converges reliably across all four system benchmarks:
\begin{itemize}
    \item \textbf{Monotonic Error Scaling:} On TFIM 10-qubit, the ground-state gap contracts monotonically from $0.4298$ at $M=100$ shots to $0.1086 \pm 0.0515$ at $M=1{,}000$ shots and $0.0324 \pm 0.0078$ at $M=10{,}000$ shots, approaching the statevector asymptotic limit ($0.0098$). On TFIM 12-qubit, the error follows an identical scaling trajectory ($0.1356$ at $1{,}000$ shots, $0.0586$ at $10{,}000$ shots).
    \item \textbf{Analytical Shot-Noise Interaction:} In finite-shot regimes, each parameter-shift gradient evaluation contains an unbiased measurement noise component $\tilde{g}_i = g_i + \xi_i$ with $\operatorname{Var}[\xi_i] = \sigma_{\text{shot}}^2 / M$. The Welford accumulator tracks the total variance $\hat{\sigma}_i^2 \approx \sigma_{\text{landscape}}^2 + \sigma_{\text{shot}}^2 / M$. Consequently, finite-shot sampling introduces an automatic additive regularization floor that scales inversely with $M$. Rather than destabilizing descent, this shot-induced variance floor inflates $F_i(t)$, dampening update magnitudes precisely when gradient signals are noisy, and gracefully shrinking as shot precision increases.
    \item \textbf{Resilience on Frustrated Spin Chains:} On the Heisenberg XXZ model, the final energy gap stabilizes at $1.7688$ ($1{,}000$ shots) and $1.4037$ ($10{,}000$ shots) at 10 qubits, which matches the noiseless statevector optimum ($1.3565 \pm 0.6480$) within statistical error.
\end{itemize}

In the budget-matched head-to-head comparison against CMA-ES (Figure~\ref{fig:shot_noise_cma_consolidated}(b) and Table~\ref{tab:shot_noise_cma_consolidated}, Part II), SIGMA-QGD statistically outperforms CMA-ES across all benchmark regimes ($p \le 0.0125$, with $\sigma_0 = 0.3$ and standard population size $\lambda = 14$). Despite CMA-ES evaluating over $4\times$ more distribution-update generations ($\sim 870$ generations vs. $200$ gradient steps for $p=30$), its rank-based selection is severely corrupted by finite-shot statistical fluctuations. In contrast, SIGMA-QGD's temporal variance accumulation and momentum filtering provide inherent robustness against shot noise.

\subsection{Dimensional Scaling Behavior}
\label{sec:scaling_results}
To assess how optimization dynamics scale with problem complexity, we analyze the transition from $n=10$ ($p=30$) to $n=12$ ($p=36$). As system size increases by two qubits, the underlying Hilbert space dimension quadruples from $2^{10} = 1,024$ to $2^{12} = 4,096$ quantum states, while the parameter manifold expands from $30$ to $36$ continuous rotation angles. 

Table~\ref{tab:dimensional_scaling} summarizes the scaling resilience across representative optimizer families on both the commuting TFIM and frustrated XXZ spin chains ($N=20$ independent random seeds, $200$ steps).

\begin{table*}[t]
\centering
\small
\caption{\textbf{Dimensional scaling behavior under Hilbert space quadrupling ($2^{10} \to 2^{12}$, $N=20$ independent seeds).} Comparison of ground-state energy gaps $|E(\theta^*) - E_0|$ and dimensional degradation ratios $\Delta E_{12\text{q}} / \Delta E_{10\text{q}}$ across key optimizer families.}
\label{tab:dimensional_scaling}
\label{tab:scaling}
\resizebox{\textwidth}{!}{%
\begin{tabular}{lcccccc}
\toprule
\multirow{2}{*}{\textbf{Optimizer Family}} & \multicolumn{3}{c}{\textbf{TFIM Spin Model}} & \multicolumn{3}{c}{\textbf{Heisenberg XXZ Spin Model}} \\
\cmidrule(lr){2-4} \cmidrule(lr){5-7}
& \textbf{10-Qubit ($p=30$)} & \textbf{12-Qubit ($p=36$)} & \textbf{Scaling Ratio ($12\text{q}/10\text{q}$)} & \textbf{10-Qubit ($p=30$)} & \textbf{12-Qubit ($p=36$)} & \textbf{Scaling Ratio ($12\text{q}/10\text{q}$)} \\
\midrule
\textbf{SIGMA-QGD v9.0 (Ours)} & $\mathbf{0.0092 \pm 0.0019}$ & $\mathbf{0.0112 \pm 0.0013}$ & $\mathbf{1.22\times}$ & $\mathbf{1.4889 \pm 0.7845}$ & $\mathbf{1.8408 \pm 0.8336}$ & $\mathbf{1.24\times}$ \\
Adam (First-Order)             & $0.1074 \pm 0.1128$          & $0.1700 \pm 0.1016$          & $1.58\times$          & $2.0897 \pm 0.9620$          & $2.1438 \pm 0.8288$          & $1.03\times$ \\
AdaBelief (Mom. Var)           & $0.0907 \pm 0.1168$          & $0.1421 \pm 0.0959$          & $1.57\times$          & $2.0853 \pm 0.9609$          & $2.0653 \pm 0.7821$          & $0.99\times$ \\
Diag-Hessian (Exact 2nd-Order) & $0.4384 \pm 0.0043$          & $0.5609 \pm 0.0075$          & $1.28\times$          & $2.1888 \pm 0.9059$          & $2.7090 \pm 0.9432$          & $1.24\times$ \\
Diag-QNG (Quantum Natural)     & $0.4068 \pm 0.0351$          & $0.4172 \pm 0.0755$          & $1.03\times$          & $1.6895 \pm 0.8629$          & $1.9602 \pm 0.7913$          & $1.16\times$ \\
COBYLA (Derivative-Free)       & $0.5194 \pm 0.0152$          & $0.6781 \pm 0.0142$          & $1.31\times$          & $3.5387 \pm 1.2916$          & $4.8212 \pm 1.7641$          & $1.36\times$ \\
SPSA (Stochastic Perturbation) & $0.5300 \pm 0.0269$          & $0.7259 \pm 0.0496$          & $1.37\times$          & $7.0031 \pm 2.4327$          & $7.3224 \pm 2.2899$          & $1.05\times$ \\
\bottomrule
\end{tabular}%
}
\end{table*}

On the TFIM model, SIGMA-QGD maintains an optimality gap an order of magnitude lower than all baselines ($\Delta E \approx 0.011$), with error increasing by only $1.22\times$ despite the $4\times$ Hilbert space expansion. In contrast, Adam and AdaBelief experience substantial degradation ($1.58\times$), resulting in final energy gaps more than $12\times$ larger than SIGMA-QGD ($0.142\text{--}0.170$ vs. $0.011$). On the frustrated XXZ model, SIGMA-QGD similarly retains superior convergence ($1.841$), scaling stably across the transition ($1.24\times$).

% ==============================================================================
% SECTION 6: DISCUSSION
% ==============================================================================
\section{Discussion}
\label{sec:discussion}

\subsection{Temporal Smoothing vs. Curvature Proxy Identity}
A key question is whether SIGMA-QGD's advantage over same-budget second-order baselines (Hess-Mag and Diag-Hessian) arises from \emph{what} is measured (gradient variance vs. instantaneous curvature) or \emph{how} it is aggregated over time (temporal accumulator vs. single-point evaluation).

Both Hess-Mag and Diag-Hessian compute the exact instantaneous parameter-shift Hessian diagonal $H_{ii}(\boldsymbol{\theta}_t)$ at the current point via Eq.~\eqref{eq:exact_hessian}. However, in non-convex quantum energy landscapes—particularly under non-commuting terms in XXZ or statistical shot noise—instantaneous second derivatives fluctuate sharply from iteration to iteration. When $H_{ii}$ transiently approaches zero near inflection points, the raw inverse preconditioner $1/(|H_{ii}| + \lambda_{\text{reg}})$ spikes, producing erratic parameter updates that destabilize convergence.

A natural question is whether augmenting instantaneous Hessian baselines with first-order momentum (analogous to quasi-Newton methods with Polyak or EMA momentum) would resolve this instability. Crucially, it cannot: if the instantaneous second derivative vanishes ($H_{ii}(\boldsymbol{\theta}_t) \approx 0$), the diagonal preconditioning factor $(|H_{ii}(\boldsymbol{\theta}_t)| + \lambda_{\text{reg}})^{-1}$ inflates to its upper regularization bound $\lambda_{\text{reg}}^{-1} = 200$ (under operational $\lambda_{\text{reg}} = 0.005$). Multiplying this catastrophic step multiplier by an accumulated momentum vector $\hat{\mathbf{m}}_t$ does not attenuate the update; rather, it injects a severe $200\times$ step-size spike along coordinate $i$, catapulting parameters out of the local basin. In contrast, SIGMA-QGD preconditions by the temporal variance $\hat{\sigma}_i^2(t)$, which aggregates gradient dispersion over the entire descent history. Because the cumulative sum of squared deviations $\sum_{\tau=1}^t (g_{c,\tau,i} - \bar{g}_{t,i})^2$ monotonically accumulates along any non-trivial trajectory, $\hat{\sigma}_i^2(t)$ never collapses to near-zero at isolated inflection points, completely avoiding the singularity spikes inherent to instantaneous Hessian inversion.

In contrast, the Welford accumulator in SIGMA-QGD computes running sample variance $\hat{\sigma}_i^2 = \operatorname{Var}_{W_t}[g_i]$ over a temporal trajectory window. This provides two distinct physical advantages:
\begin{itemize}
    \item \textbf{Low-Pass Curvature Regularization:} Temporal accumulation via Welford's algorithm computes the cumulative sample variance over the trajectory history, damping high-frequency gradient fluctuations without requiring extra quantum circuits.
    \item \textbf{Path-Integrated Sensitivity vs. Static Point Curvature:} By Proposition~\ref{prop:variance_curvature}, $\hat{\sigma}_i^2 \approx H_{ii}^2 \operatorname{Var}_{W_t}[\theta_i]$. Rather than measuring static local curvature at an isolated point $\boldsymbol{\theta}_t$, trajectory variance measures dynamical sensitivity along the \emph{realized descent path}. Coordinates experiencing large gradient swings across recent steps receive conservative, stabilizing scaling, tailoring step sizes directly to the geometry of the traversed ravine.
\end{itemize}
Temporal smoothing is therefore not a confounding artifact, but an essential structural property of variance-based preconditioning that naturally regularizes quantum landscapes at zero additional circuit overhead. This observation aligns directly with recent theoretical treatments of curvature regularization and metric damping in quantum variational optimization \cite{stokes2020quantum,wierichs2020avoiding}, which establish that unregularized second-order steps severely destabilize optimization in non-convex quantum landscapes unless dampened along ill-conditioned directions.

\subsection{Ansatz Expressivity and Landscape Frustration}
On the non-commuting XXZ Hamiltonian, diagonal QNG and adaptive first-order methods achieve comparable final energy gaps to SIGMA-QGD ($\Delta E \approx 1.5\text{--}1.8$). This residual optimality gap highlights an important distinction between optimization performance and ansatz expressivity:
\begin{itemize}
    \item \textbf{Ansatz Expressivity vs. Optimization Limits:} The 1D Heisenberg XXZ chain exhibits gapless critical ground states and non-local entanglement that cannot be represented exactly by a shallow $r=3$ layer hardware-efficient $RY$-CNOT circuit. Consequently, the $\sim 1.5$ gap reflects the representational bound of the variational state manifold itself, rather than premature optimization failure.
    \item \textbf{Shallow NISQ Regimes vs. Barren Plateaus:} While deeper ans\"atze ($r \ge 6$) would expand expressivity and reduce the ground-state gap, they inevitably suffer from noise-induced barren plateaus on physical NISQ hardware. By benchmarking in this expressivity-constrained, highly frustrated regime, our experiments evaluate whether optimizers can find the global variational optimum of a realistic NISQ circuit without being trapped by spurious saddle points or directional thrashing.
    \item \textbf{Boundary of Diagonal Approximations:} The empirical results also delineate the natural boundary of diagonal curvature preconditioning: in quantum models dominated by dense, non-local Pauli exchange interactions, diagonal projections capture only a fraction of the total landscape curvature, explaining why QNG and SIGMA-QGD converge to similar asymptotic energy plateaus.
\end{itemize}

\subsection{Limitations}
We explicitly note the following methodological bounds:
\begin{itemize}
    \item \textbf{Ansatz Specificity:} Benchmarks were performed on hardware-efficient $RY$-CNOT ans\"atze; problem-tailored ans\"atze (such as unitary coupled-cluster circuits) require further investigation.
    \item \textbf{Stationary Point vs. Global Basin Guarantees:} While Theorem~\ref{thm:convergence} establishes an ergodic $\mathcal{O}(1/T)$ rate to first-order stationary points ($\min_{1 \le t \le T} \|\nabla E(\boldsymbol{\theta}_{t-1})\|_2^2 \le \epsilon$) under uniform margin conditions, global convergence guarantees to the true ground-state energy basin remain an open theoretical challenge on highly frustrated, non-convex quantum manifolds.
\end{itemize}

% ==============================================================================
% SECTION 7: CONCLUSION
% ==============================================================================
\section{Conclusion}
\label{sec:conclusion}

SIGMA-QGD introduces a zero-extra-circuit, variance-aware curvature proxy that bridges the gap between first-order gradient methods and metric-aware quantum natural gradients. By extracting curvature information from the temporal statistics of parameter-shift evaluations, SIGMA-QGD achieves consistent convergence gains across spin models and system sizes without increasing quantum circuit requirements beyond the standard $2p+1$ evaluations.

Future work will focus on:
\begin{itemize}
    \item Fusing the trajectory-variance proxy with exact analytic Hessian diagonals to capture both historical and instantaneous curvature signals;
    \item Validating performance on physical NISQ quantum processors subject to device noise;
    \item Extending the framework to chemistry Hamiltonians and unitary coupled-cluster (UCC) ans\"atze.
\end{itemize}

% ==============================================================================
% CODE AND DATA AVAILABILITY
% ==============================================================================
\section*{Code and Data Availability}
The complete implementation, YAML experimental configurations, raw per-seed benchmark CSV datasets, and figure generation scripts are publicly available on GitHub at:\\
\url{https://github.com/Arjun-E-Naik/SIGMA-QGD}.

% ==============================================================================
% REFERENCES
% ==============================================================================
\begin{thebibliography}{99}

\bibitem{peruzzo2014variational}
A.~Peruzzo, P.~McClean, P.~Shadbolt, M.-H.~Yung, X.-Q.~Zhou, P.~J.~Love, A.~Aspuru-Guzik, and J.~L.~O'Brien,
\newblock ``A variational eigenvalue solver on a photonic quantum processor,''
\newblock \emph{Nature Communications}, vol.~5, p.~4213, 2014.

\bibitem{mcclean2016theory}
J.~R.~McClean, J.~Romero, R.~Babbush, and A.~Aspuru-Guzik,
\newblock ``The theory of variational hybrid quantum-classical algorithms,''
\newblock \emph{New Journal of Physics}, vol.~18, no.~2, p.~023023, 2016.

\bibitem{mitarai2018quantum}
K.~Mitarai, M.~Negoro, M.~Kitagawa, and K.~Fujii,
\newblock ``Quantum circuit learning,''
\newblock \emph{Physical Review A}, vol.~98, no.~3, p.~032309, 2018.

\bibitem{schuld2019evaluating}
M.~Schuld, V.~Bergholm, C.~Gogolin, J.~Izaac, and N.~Killoran,
\newblock ``Evaluating analytic gradients on quantum hardware,''
\newblock \emph{Physical Review A}, vol.~99, no.~3, p.~032331, 2019.

\bibitem{mcclean2018barren}
J.~R.~McClean, S.~Boixo, V.~N.~Smelyanskiy, R.~Babbush, and H.~Neven,
\newblock ``Barren plateaus in quantum neural network training landscapes,''
\newblock \emph{Nature Communications}, vol.~9, p.~4812, 2018.

\bibitem{cerezo2021cost}
M.~Cerezo, A.~Sone, T.~Volkoff, L.~Cincio, and P.~J.~Coles,
\newblock ``Cost function dependent barren plateaus in shallow parametrized quantum circuits,''
\newblock \emph{Nature Communications}, vol.~12, p.~1791, 2021.

\bibitem{kingma2014adam}
D.~P.~Kingma and J.~Ba,
\newblock ``Adam: A method for stochastic optimization,''
\newblock \emph{arXiv:1412.6980}, 2014.

\bibitem{tieleman2012lecture}
T.~Tieleman and G.~Hinton,
\newblock ``Lecture 6.5---RMSProp: Divide the gradient by a running average of its recent magnitude,''
\newblock \emph{COURSERA: Neural Networks for Machine Learning}, 2012.

\bibitem{zaheer2018adaptive}
M.~Zaheer, S.~Reddi, D.~Sachan, S.~Kale, and S.~Kumar,
\newblock ``Adaptive methods for nonconvex optimization,''
\newblock \emph{Advances in Neural Information Processing Systems (NeurIPS)}, vol.~31, 2018.

\bibitem{zhuang2020adabelief}
J.~Zhuang, T.~Tang, Y.~Ding, S.~C.~Tatikonda, N.~Dvornek, X.~Papademetris, and J.~Duncan,
\newblock ``AdaBelief optimizer: Adapting stepsizes by the belief in observed gradients,''
\newblock \emph{Advances in Neural Information Processing Systems (NeurIPS)}, vol.~33, 2020.

\bibitem{stokes2020quantum}
J.~Stokes, J.~Izaac, N.~Killoran, and G.~Carleo,
\newblock ``Quantum natural gradient,''
\newblock \emph{Quantum}, vol.~4, p.~269, 2020.

\bibitem{wierichs2020avoiding}
D.~Wierichs, C.~Gogolin, and M.~Kastoryano,
\newblock ``Avoiding local minima in variational quantum eigensolvers with the natural gradient,''
\newblock \emph{Physical Review Research}, vol.~2, no.~4, p.~043246, 2020.

\bibitem{spall1992multivariate}
J.~C.~Spall,
\newblock ``Multivariate stochastic approximation using a simultaneous perturbation gradient approximation,''
\newblock \emph{IEEE Transactions on Automatic Control}, vol.~37, no.~3, pp.~332--341, 1992.

\bibitem{novak2025benchmarking}
V.~Nov{\'a}k, I.~Zelinka, and V.~Sn{\'a}{\v{s}}el,
\newblock ``Optimization strategies for variational quantum algorithms in noisy landscapes,''
\newblock \emph{arXiv:2506.01715}, 2025.

\bibitem{shi2025hamiltonian}
C.~Shi and H.~Wang,
\newblock ``Efficient Hamiltonian-aware quantum natural gradient descent for variational quantum eigensolvers,''
\newblock \emph{arXiv:2511.14511}, 2025.

\bibitem{halla2025stein}
M.~Halla,
\newblock ``Estimation of quantum Fisher information via Stein's identity in variational quantum algorithms,''
\newblock \emph{arXiv:2502.17231}, 2025.

\bibitem{wierichs2022general}
D.~Wierichs, J.~Izaac, C.~Wang, and C.~Y.-Y.~Lin,
\newblock ``General parameter-shift rules for quantum gradients,''
\newblock \emph{Quantum}, vol.~6, p.~677, 2022.

\bibitem{powell1994direct}
M.~J.~D.~Powell,
\newblock ``A direct search optimization method that models the objective and constraint functions by linear interpolation,''
\newblock \emph{Advances in Optimization and Numerical Analysis}, pp.~51--67, Springer, 1994.

\bibitem{kandala2017hardware}
A.~Kandala, A.~Mezzacapo, K.~Temme, M.~Takita, M.~Brink, J.~M.~Chow, and J.~M.~Gambetta,
\newblock ``Hardware-efficient variational quantum eigensolver for small molecules and quantum magnets,''
\newblock \emph{Nature}, vol.~549, pp.~242--246, 2017.

\bibitem{welford1962note}
B.~P.~Welford,
\newblock ``Note on a method for calculating corrected sums of squares and products,''
\newblock \emph{Technometrics}, vol.~4, no.~3, pp.~419--420, 1962.

\bibitem{gacon2021simultaneous}
J.~Gacon, C.~Zoufal, G.~Carleo, and S.~Woerner,
\newblock ``Simultaneous perturbation stochastic approximation of the quantum Fisher information,''
\newblock \emph{Quantum}, vol.~5, p.~567, 2021.

\bibitem{spall1997accelerated}
J.~C.~Spall,
\newblock ``Accelerated second-order stochastic optimization using only function measurements,''
\newblock \emph{Proceedings of the 36th IEEE Conference on Decision and Control (CDC)}, vol.~2, pp.~1417--1424, 1997.

\bibitem{spall1998implementation}
J.~C.~Spall,
\newblock ``Implementation of the simultaneous perturbation algorithm for stochastic optimization,''
\newblock \emph{IEEE Transactions on Aerospace and Electronic Systems}, vol.~34, no.~3, pp.~817--823, 1998.

\bibitem{spall2000adaptive}
J.~C.~Spall,
\newblock ``Adaptive simultaneous perturbation stochastic approximation for simulated systems,''
\newblock \emph{IEEE Transactions on Aerospace and Electronic Systems}, vol.~36, no.~2, pp.~439--448, 2000.

\bibitem{qnspsa_psr_2024}
D.-T.~Le, V.-L.~Nguyen, H.~C.~Nguyen, Q.-H.~Nguyen, and V.-D.~Nguyen,
\newblock ``Variational quantum eigensolver: A comparative analysis of classical and quantum optimizer methods,''
\newblock \emph{arXiv:2412.19176}, 2024.

\bibitem{hessian_resolvability_2026}
Y.~Huang and Y.~Wang,
\newblock ``Analysis of Hessian scaling for local and global costs in variational quantum algorithms,''
\newblock \emph{arXiv:2602.00783}, 2026.

\bibitem{ghadimi2013stochastic}
S.~Ghadimi and G.~Lan,
\newblock ``Stochastic first- and zeroth-order methods for nonconvex stochastic programming,''
\newblock \emph{SIAM Journal on Optimization}, vol.~23, no.~4, pp.~2341--2368, 2013.

\end{thebibliography}

\end{document}
